\documentclass[12pt,a4paper]{article}
\usepackage[utf8]{inputenc}
\usepackage{amsmath}
\usepackage{graphicx}
\usepackage[hidelinks]{hyperref}
\usepackage{geometry}
\geometry{a4paper, margin=1in}
\usepackage{listings}
\usepackage{xcolor}
\usepackage{booktabs}
\usepackage{caption}

\definecolor{codegreen}{rgb}{0,0.6,0}
\definecolor{codegray}{rgb}{0.5,0.5,0.5}
\definecolor{codepurple}{rgb}{0.58,0,0.82}
\definecolor{backcolour}{rgb}{0.95,0.95,0.92}

\lstdefinestyle{mystyle}{
    backgroundcolor=\color{backcolour},   
    commentstyle=\color{codegreen},
    keywordstyle=\color{magenta},
    numberstyle=\tiny\color{codegray},
    stringstyle=\color{codepurple},
    basicstyle=\ttfamily\footnotesize,
    breakatwhitespace=false,         
    breaklines=true,                 
    captionpos=b,                    
    keepspaces=true,                 
    numbers=left,                    
    numbersep=5pt,                  
    showspaces=false,                
    showstringspaces=false,
    showtabs=false,                  
    tabsize=2
}

\lstset{style=mystyle}

\title{
    \vspace{-2cm}
    \begin{center}
        \Large \textbf{German University in Cairo (GUC)} \\
        \large Digital System Design - 6th Semester \\
    \end{center}
    \vspace{2cm}
    \textbf{FPGA-Accelerated CNN Inference Engine} \\
    \large A Custom VHDL Neural Network Accelerator on Artix-7
}
\author{
    \textbf{Team: 7 Segment Displays} \\[0.5cm]
    \normalsize
    \begin{tabular}{c}
        Muhammad Walid 61-34156 T4 \\
        Morkos Shenouda [ID] [TUT] \\
        Youssef Tarek [ID] [TUT] \\
        Ziad Saad [ID] [TUT] \\
        Amr Ahmed [ID] [TUT] \\
        David Hany [ID] [TUT] \\
        Ziad Ahmed [ID] [TUT]
    \end{tabular}
}
\date{\vspace{1cm}
    \textbf{Professor:} Dr. Mohamed Abdel Ghany Ahmed Salem \\
    \textbf{TA:} Eng. Mohamed Ahmed Awaad Ibrahim \\
    \today
}

\begin{document}

\maketitle
\newpage
\tableofsections
\tableofcontents
\newpage

\section{Abstract}
This project presents the design, implementation, and verification of a custom FPGA-accelerated Convolutional Neural Network (CNN) inference engine. The system is capable of classifying 32x32 grayscale images into three categories (Cat, Dog, Human) using a fully custom VHDL datapath. It avoids High-Level Synthesis (HLS) and vendor IP cores for the CNN logic, achieving near cycle-parity (0.90x) with a 5.2 GHz Intel i7-14700HX processor while operating at a 52x lower clock frequency (100 MHz) and consuming only 0.269W. The project evolved in two phases: Phase 1 utilized an SPI interface for communication with a Raspberry Pi, and Phase 2 upgraded to a 100 Mbps Ethernet (RMII) interface to eliminate I/O bottlenecks and enable rigorous desktop benchmarking.

\section{Introduction}
General-purpose CPUs execute sequential instructions, limiting their efficiency in parallel workloads like CNN inference. This project aims to exploit the spatial parallelism of FPGAs to accelerate a 5-layer CNN model. 

The custom datapath leverages the Artix-7 FPGA's DSP48E1 slices and Block RAMs to perform concurrent Multiply-Accumulate (MAC) operations. By adopting a Q8.8 16-bit fixed-point arithmetic format, the design maintains bit-exact parity with a software Golden Model.

\section{System Architecture}

\subsection{CNN Model Overview}
The CNN was trained in Python using TensorFlow/Keras on a custom dataset over 50 epochs. The architecture consists of:
\begin{itemize}
    \item Input: 32x32x1 (Padded to 34x34 for hardware)
    \item Conv1 + ReLU: 3x3, 16 filters
    \item MaxPool1: 2x2
    \item Conv2 + ReLU: 3x3, 32 filters
    \item MaxPool2: 2x2
    \item Conv3 + ReLU: 3x3, 32 filters
    \item MaxPool3: 2x2
    \item Flatten
    \item Dense1 + ReLU: 512 to 32
    \item Dense2 (Logits): 32 to 3
\end{itemize}

\subsection{Hardware Datapath}
The hardware datapath is composed of 13 custom VHDL modules. 
The convolution layers utilize line buffers and a sliding window mechanism to feed data into a DSP48E1-based binary adder tree. The Dense layers feature parallel MAC engines to rapidly compute matrix multiplications. Inter-module communication is governed by a robust AXI-Stream-like valid/ready handshake, ensuring backpressure support and data integrity.

\section{Implementation Phases}

\subsection{Phase 1: SPI Communication}
In the first phase, a Raspberry Pi 5 served as the host, transmitting image pixels and receiving logits via an 8 MHz SPI bus. While functionally correct and significantly faster than the Pi's ARM CPU in pure compute (15x fewer cycles), the system was I/O bound. The SPI transfer accounted for 69\% of the total latency, masking the true performance of the FPGA accelerator.

\subsection{Phase 2: Ethernet Upgrade}
To overcome the I/O bottleneck, Phase 2 integrated a 100 Mbps RMII Ethernet stack. An open-source RMII MAC was paired with a custom \texttt{eth\_cnn\_bridge} module, incorporating RX and TX FIFOs. The bridge automatically strips the 42-byte OSI headers (Ethernet + IP + UDP) and streams the 1,156-byte image payload directly into the CNN core. This upgrade provided 12.5x more bandwidth, allowing direct performance comparisons against a high-end desktop CPU.

\section{Verification and Benchmarks}

\subsection{Verification Methodology}
A Python-based integer-only Golden Model (\texttt{golden\_model.py}) was developed to perfectly emulate the FPGA's Q8.8 fixed-point arithmetic, including DSP48 MAC behavior and ReLU activations. 100 test images were evaluated on both the hardware and the Golden Model, achieving 100\% bit-exact parity.

\subsection{Performance Metrics}
The Ethernet-enabled FPGA was benchmarked against an Intel i7-14700HX running TensorFlow Lite:
\begin{itemize}
    \item \textbf{Intel i7 (5.2 GHz):} 92 $\mu$s latency, 477,880 cycles, ~45W TDP.
    \item \textbf{Artix-7 (100 MHz):} 5.3 ms latency, 530,759 cycles, 0.269W TDP.
    \item \textbf{Architectural Efficiency:} The FPGA achieved a 0.90x cycle-parity with the i7, demonstrating that the custom spatial architecture requires only 11\% more cycles than a superscalar x86 processor.
    \item \textbf{Power Efficiency:} The FPGA is 167x more power-efficient per inference.
\end{itemize}

\subsection{Resource Utilization}
\begin{itemize}
    \item \textbf{LUTs:} 3,919 (6.18\%)
    \item \textbf{Flip-Flops:} 3,634 (2.87\%)
    \item \textbf{BRAM:} 15.5 (11.48\%)
    \item \textbf{DSP48E1:} 46 (19.17\%)
    \item \textbf{WNS:} +7.042 ns
\end{itemize}

\section{Conclusion and Future Work}
The project successfully delivered a fully custom, high-efficiency CNN accelerator. The core logic proved interface-agnostic, seamlessly transitioning from SPI to Ethernet while maintaining bit-exact accuracy. 
Future work includes migrating the design to a Zynq-7000 SoC to leverage shared AXI memory, eliminating I/O transit overhead completely and potentially eclipsing desktop CPU latency in absolute time.

\newpage
\appendix
\section{VHDL Source Code}

\subsection{cnn_config.vhd}
\begin{lstlisting}[language=VHDL]
library IEEE;
use IEEE.STD_LOGIC_1164.ALL;
use IEEE.NUMERIC_STD.ALL;

package cnn_config is
    -- Image dimensions
    constant IMG_W            : integer := 32;
    constant IMG_H            : integer := 32;
    constant PADDED_IMG_W     : integer := 34; -- Pi-side padding
    constant PADDED_IMG_H     : integer := 34;
    constant IMG_SIZE         : integer := 1156; -- 34*34

    -- Layer output dimensions (after valid conv, before pool)
    constant CONV1_OUT_W      : integer := 32; -- 34-2
    constant POOL1_OUT_W      : integer := 16; -- 32/2
    constant CONV2_IN_W       : integer := 18; -- 16+2 (Padder1)
    constant CONV2_OUT_W      : integer := 16; -- 18-2
    constant POOL2_OUT_W      : integer := 8;  -- 16/2
    constant CONV3_IN_W       : integer := 10; -- 8+2 (Padder2)
    constant CONV3_OUT_W      : integer := 8;  -- 10-2
    constant POOL3_OUT_W      : integer := 4;  -- 8/2
    constant FLATTEN_SIZE     : integer := 512; -- 4*4*32

    -- Channel counts
    constant CONV1_IN_CH      : integer := 1;
    constant CONV1_OUT_CH     : integer := 16;
    constant CONV2_IN_CH      : integer := 16;
    constant CONV2_OUT_CH     : integer := 32;
    constant CONV3_IN_CH      : integer := 32;
    constant CONV3_OUT_CH     : integer := 32;
    constant DENSE1_IN        : integer := 512;
    constant DENSE1_OUT       : integer := 32;
    constant DENSE2_IN        : integer := 32;
    constant DENSE2_OUT       : integer := 3;

    -- Arithmetic
    constant DATA_WIDTH       : integer := 16; -- Q8.8
    constant ACCUM_WIDTH      : integer := 48; -- DSP48 P register
    constant FRAC_BITS        : integer := 8;  -- fractional bits

    -- SPI
    constant SPI_DATA_WIDTH   : integer := 16;

    -- Parallelism (set to 2 for recommended configuration)
    constant CONV_OC_PARALLEL : integer := 2; -- output channels computed in parallel
    constant DENSE_PARALLEL   : integer := 8; -- MACs in dense engine

    -- Derived: number of classes
    constant NUM_CLASSES      : integer := 3;
end package cnn_config;
\end{lstlisting}

\subsection{system_top.vhd}
\begin{lstlisting}[language=VHDL]
library IEEE;
use IEEE.STD_LOGIC_1164.ALL;
use IEEE.NUMERIC_STD.ALL;
use work.cnn_config.all;

entity system_top is
    port (
        clk_100mhz : in  std_logic;
        reset_n    : in  std_logic; 
        
        -- SPI Interface
        spi_cs_n   : in  std_logic;
        spi_sclk   : in  std_logic;
        spi_mosi   : in  std_logic;
        spi_miso   : out std_logic;
        
        -- Hardware Interrupt
        data_ready : out std_logic;
        
        -- NEW: Hardware Debugger LEDs
        led        : out std_logic_vector(15 downto 0)
    );
end entity;

architecture struct of system_top is

    signal sys_reset      : std_logic;
    signal soft_reset     : std_logic;
    signal combined_reset : std_logic;

    -- SPI Controller Signals
    signal spi_rx_data    : std_logic_vector(15 downto 0);
    signal spi_rx_valid   : std_logic;
    signal spi_tx_data    : std_logic_vector(15 downto 0);
    signal spi_tx_load    : std_logic;

    -- Synchronization (The Metastability Fix)
    signal sclk_sync      : std_logic_vector(2 downto 0) := "000";
    signal cs_sync        : std_logic_vector(1 downto 0) := "11";

    -- FIFO and CNN Signals
    signal rx_fifo_wr_en, rx_fifo_rd_en, rx_fifo_empty, rx_fifo_full : std_logic;
    signal rx_fifo_data   : std_logic_vector(15 downto 0);
    signal tx_fifo_wr_en, tx_fifo_rd_en, tx_fifo_empty, tx_fifo_full : std_logic;
    signal tx_fifo_data   : std_logic_vector(15 downto 0);
    signal cnn_in_valid, cnn_in_ready, cnn_out_valid, cnn_out_ready : std_logic;
    signal cnn_out_data   : std_logic_vector(15 downto 0);

    -- FSM Definition
    type system_state_type is (IDLE, CAPTURE_AND_INFERENCE, PRE_LOAD, SEND_RESULTS, WAIT_FOR_SAFE_LOAD);
    signal state : system_state_type := IDLE;

    signal pixel_count : integer range 0 to 2047 := 0;
    signal logit_count : integer range 0 to 15 := 0;
    signal read_count  : integer range 0 to 15 := 0;

begin

    sys_reset <= not reset_n;
    combined_reset <= sys_reset or soft_reset;

    -- Synchronizers for Top-Level FSM (fixes metastability!)
    process(clk_100mhz)
    begin
        if rising_edge(clk_100mhz) then
            sclk_sync <= sclk_sync(1 downto 0) & spi_sclk;
            cs_sync   <= cs_sync(0) & spi_cs_n;
        end if;
    end process;

    u_spi_ctrl: entity work.spi_controller
        generic map(G_TX_WIDTH => 16, G_RX_WIDTH => 16)
        port map(
            clk_100mhz => clk_100mhz, reset => sys_reset,
            spi_cs_n => spi_cs_n, spi_sclk => spi_sclk, spi_mosi => spi_mosi, spi_miso => spi_miso,
            tx_data => spi_tx_data, tx_load => spi_tx_load, rx_data => spi_rx_data, rx_valid => spi_rx_valid
        );

    u_rx_fifo: entity work.sync_fifo
        generic map(DATA_WIDTH => 16, ADDR_WIDTH => 11)
        port map(clk => clk_100mhz, reset => combined_reset, wr_en => rx_fifo_wr_en, wr_data => spi_rx_data, rd_en => rx_fifo_rd_en, rd_data => rx_fifo_data, empty => rx_fifo_empty, full => rx_fifo_full);

    u_cnn_core: entity work.cnn_datapath_top
        port map(clk => clk_100mhz, reset => combined_reset, in_data => rx_fifo_data, in_valid => cnn_in_valid, in_ready => cnn_in_ready, out_data => cnn_out_data, out_valid => cnn_out_valid, out_ready => cnn_out_ready);

    u_tx_fifo: entity work.sync_fifo
        generic map(DATA_WIDTH => 16, ADDR_WIDTH => 4)
        port map(clk => clk_100mhz, reset => combined_reset, wr_en => tx_fifo_wr_en, wr_data => cnn_out_data, rd_en => tx_fifo_rd_en, rd_data => tx_fifo_data, empty => tx_fifo_empty, full => tx_fifo_full);

    -- Routing
    rx_fifo_wr_en <= spi_rx_valid when state = CAPTURE_AND_INFERENCE else '0';
    cnn_in_valid  <= not rx_fifo_empty;
    rx_fifo_rd_en <= cnn_in_ready and not rx_fifo_empty;
    tx_fifo_wr_en <= cnn_out_valid;
    cnn_out_ready <= not tx_fifo_full; 
    spi_tx_data   <= tx_fifo_data;

    -- Master FSM
    process(clk_100mhz)
    begin
        if rising_edge(clk_100mhz) then
            if sys_reset = '1' then
                state <= IDLE;
                data_ready <= '0';
                soft_reset <= '1';
                spi_tx_load <= '0';
                tx_fifo_rd_en <= '0';
                pixel_count <= 0;
                logit_count <= 0;
                read_count <= 0;
            else
                -- Default strobes
                spi_tx_load <= '0';
                tx_fifo_rd_en <= '0';
                soft_reset <= '0';

                case state is
                    when IDLE =>
                        data_ready <= '0';
                        pixel_count <= 0;
                        logit_count <= 0;
                        read_count <= 0;
                        soft_reset <= '1'; -- Hold CNN in reset until Pi starts talking
                        
                        -- Start ONLY when synchronized CS goes low
                        if cs_sync(1) = '0' then
                            state <= CAPTURE_AND_INFERENCE;
                            soft_reset <= '0';
                        end if;

                    when CAPTURE_AND_INFERENCE =>
                        if spi_rx_valid = '1' then 
                            pixel_count <= pixel_count + 1;
                        end if;
                        
                        if cnn_out_valid = '1' then
                            logit_count <= logit_count + 1;
                            if logit_count = NUM_CLASSES - 1 then 
                                state <= PRE_LOAD; 
                            end if;
                        end if;

                        -- Emergency abort if Pi raises CS_N prematurely
                        if cs_sync(1) = '1' and logit_count < NUM_CLASSES - 1 then
                            state <= IDLE;
                        end if;

                    when PRE_LOAD =>
                        spi_tx_load   <= '1';
                        tx_fifo_rd_en <= '1'; 
                        state <= SEND_RESULTS;

                    when SEND_RESULTS =>
                        data_ready <= '1';
                        
                        if spi_rx_valid = '1' then
                            read_count <= read_count + 1;
                            state <= WAIT_FOR_SAFE_LOAD;
                        end if;

                        -- If Pi raises CS_N, the transaction is over!
                        if cs_sync(1) = '1' then
                            state <= IDLE;
                        end if;

                    when WAIT_FOR_SAFE_LOAD =>
                        data_ready <= '1'; -- Keep it high so Pi doesn't drop out
                        if sclk_sync(1) = '0' then
                            tx_fifo_rd_en <= '1';
                            spi_tx_load   <= '1'; 
                            state <= SEND_RESULTS;
                        end if;

                        if cs_sync(1) = '1' then
                            state <= IDLE;
                        end if;

                    when others => 
                        state <= IDLE;
                end case;
            end if;
        end if;
    end process;

    -- ==========================================
    -- HARDWARE LED DEBUGGER
    -- ==========================================
    led(15 downto 12) <= "0000" when state = IDLE else
                         "0001" when state = CAPTURE_AND_INFERENCE else
                         "0010" when state = PRE_LOAD else
                         "0011" when state = SEND_RESULTS else
                         "0100" when state = WAIT_FOR_SAFE_LOAD else 
                         "1111";
                         
    -- LED 11: Flashes if the CNN is stuck waiting for pixels (RX FIFO empty)
    led(11) <= rx_fifo_empty;
    
    -- LEDs 10-0: Binary representation of the Pixel Count
    led(10 downto 0) <= std_logic_vector(to_unsigned(pixel_count, 11));

end architecture;
\end{lstlisting}

\subsection{cnn_datapath_top.vhd}
\begin{lstlisting}[language=VHDL]
library IEEE;
use IEEE.STD_LOGIC_1164.ALL;
use IEEE.NUMERIC_STD.ALL;
use work.cnn_config.all;

entity cnn_datapath_top is
    generic (
        -- Global Dimensions
        G_IMG_W         : integer := 34;
        G_IMG_H         : integer := 34;

        -- Path prefix relative to the Vivado synth_1 run directory:
        -- ../../ goes from <project>.runs/synth_1 to the project root.
        
        -- Stage 1: Conv1 (34x34 -> 32x32)
        C1_W_FILE       : string  := "../../FPGA_Accelerated_CNN_Inference.srcs/sources_1/imports/export_q88/conv1_weights.mem";
        C1_B_FILE       : string  := "../../FPGA_Accelerated_CNN_Inference.srcs/sources_1/imports/export_q88/conv1_biases.mem";
        C1_OUT_CH       : integer := 16;

        -- Stage 2: Conv2 (18x18 -> 16x16)
        C2_W_FILE       : string  := "../../FPGA_Accelerated_CNN_Inference.srcs/sources_1/imports/export_q88/conv2_weights.mem";
        C2_B_FILE       : string  := "../../FPGA_Accelerated_CNN_Inference.srcs/sources_1/imports/export_q88/conv2_biases.mem";
        C2_OUT_CH       : integer := 32;

        -- Stage 3: Conv3 (10x10 -> 8x8)
        C3_W_FILE       : string  := "../../FPGA_Accelerated_CNN_Inference.srcs/sources_1/imports/export_q88/conv3_weights.mem";
        C3_B_FILE       : string  := "../../FPGA_Accelerated_CNN_Inference.srcs/sources_1/imports/export_q88/conv3_biases.mem";
        C3_OUT_CH       : integer := 32;

        -- Stage 4: Dense1 (512 -> 32)
        D1_W_FILE       : string  := "../../FPGA_Accelerated_CNN_Inference.srcs/sources_1/imports/export_q88/dense1_weights.mem";
        D1_B_FILE       : string  := "../../FPGA_Accelerated_CNN_Inference.srcs/sources_1/imports/export_q88/dense1_biases.mem";
        D1_IN_SIZE      : integer := 512;
        D1_OUT_SIZE     : integer := 32;

        -- Stage 5: Dense2 (32 -> 8)
        D2_W_FILE       : string  := "../../FPGA_Accelerated_CNN_Inference.srcs/sources_1/imports/export_q88/dense2_weights.mem";
        D2_B_FILE       : string  := "../../FPGA_Accelerated_CNN_Inference.srcs/sources_1/imports/export_q88/dense2_biases.mem";
        D2_IN_SIZE      : integer := 32;
        D2_OUT_SIZE     : integer := 8
    );
    port (
        clk       : in  std_logic;
        reset     : in  std_logic;
        
        -- AXI-Stream Input
        in_data   : in  std_logic_vector(15 downto 0);
        in_valid  : in  std_logic;
        in_ready  : out std_logic;
        
        -- AXI-Stream Output
        out_data  : out std_logic_vector(15 downto 0);
        out_valid : out std_logic;
        out_ready : in  std_logic
    );
end entity;

architecture struct of cnn_datapath_top is

    -- Interconnect signals (Data, Valid, Ready)
    signal s1_data, s2_data, s3_data, s4_data, s5_data   : std_logic_vector(15 downto 0);
    signal s1_valid, s2_valid, s3_valid, s4_valid, s5_valid : std_logic;
    signal s1_ready, s2_ready, s3_ready, s4_ready, s5_ready : std_logic;

    signal s6_data, s7_data, s8_data, s9_data, s10_data  : std_logic_vector(15 downto 0);
    signal s6_valid, s7_valid, s8_valid, s9_valid, s10_valid : std_logic;
    signal s6_ready, s7_ready, s8_ready, s9_ready, s10_ready : std_logic;

    signal flatten_done_open : std_logic;

begin

    -- =========================================================
    -- STAGE 1: Feature Extraction
    -- =========================================================
    u_conv1: entity work.conv_layer
        generic map(
            NUM_IN_CH   => 1,
            NUM_OUT_CH  => C1_OUT_CH,
            IMG_W       => G_IMG_W,
            IMG_H       => G_IMG_H,
            OC_PARALLEL => 1,
            WEIGHT_FILE => C1_W_FILE,
            BIAS_FILE   => C1_B_FILE
        )
        port map(
            clk => clk, reset => reset,
            in_data => in_data, in_valid => in_valid, in_ready => in_ready,
            out_data => s1_data, out_valid => s1_valid, out_ready => s1_ready
        );

    u_pool1: entity work.pool_layer
        generic map(IMG_W => G_IMG_W-2, IMG_H => G_IMG_H-2, CHANNELS => C1_OUT_CH)
        port map(clk => clk, reset => reset, in_data => s1_data, in_valid => s1_valid, in_ready => s1_ready, out_data => s2_data, out_valid => s2_valid, out_ready => s2_ready);

    u_pad1: entity work.padder
        generic map(IMG_W => (G_IMG_W-2)/2, IMG_H => (G_IMG_H-2)/2, CHANNELS => C1_OUT_CH)
        port map(clk => clk, reset => reset, in_data => s2_data, in_valid => s2_valid, in_ready => s2_ready, out_data => s3_data, out_valid => s3_valid, out_ready => s3_ready);

    -- =========================================================
    -- STAGE 2: Mid-Level Features
    -- =========================================================
    u_conv2: entity work.conv_layer
        generic map(
            NUM_IN_CH   => C1_OUT_CH,
            NUM_OUT_CH  => C2_OUT_CH,
            IMG_W       => ((G_IMG_W-2)/2)+2,
            IMG_H       => ((G_IMG_H-2)/2)+2,
            OC_PARALLEL => 1,
            WEIGHT_FILE => C2_W_FILE,
            BIAS_FILE   => C2_B_FILE
        )
        port map(
            clk => clk, reset => reset,
            in_data => s3_data, in_valid => s3_valid, in_ready => s3_ready,
            out_data => s4_data, out_valid => s4_valid, out_ready => s4_ready
        );

    u_pool2: entity work.pool_layer
        generic map(IMG_W => (G_IMG_W-2)/2, IMG_H => (G_IMG_H-2)/2, CHANNELS => C2_OUT_CH)
        port map(clk => clk, reset => reset, in_data => s4_data, in_valid => s4_valid, in_ready => s4_ready, out_data => s5_data, out_valid => s5_valid, out_ready => s5_ready);

    u_pad2: entity work.padder
        generic map(IMG_W => (G_IMG_W-2)/4, IMG_H => (G_IMG_H-2)/4, CHANNELS => C2_OUT_CH)
        port map(clk => clk, reset => reset, in_data => s5_data, in_valid => s5_valid, in_ready => s5_ready, out_data => s6_data, out_valid => s6_valid, out_ready => s6_ready);

    -- =========================================================
    -- STAGE 3: High-Level Features
    -- =========================================================
    u_conv3: entity work.conv_layer
        generic map(
            NUM_IN_CH   => C2_OUT_CH,
            NUM_OUT_CH  => C3_OUT_CH,
            IMG_W       => ((G_IMG_W-2)/4)+2,
            IMG_H       => ((G_IMG_H-2)/4)+2,
            OC_PARALLEL => 1,
            WEIGHT_FILE => C3_W_FILE,
            BIAS_FILE   => C3_B_FILE
        )
        port map(
            clk => clk, reset => reset,
            in_data => s6_data, in_valid => s6_valid, in_ready => s6_ready,
            out_data => s7_data, out_valid => s7_valid, out_ready => s7_ready
        );

    u_pool3: entity work.pool_layer
        generic map(IMG_W => (G_IMG_W-2)/4, IMG_H => (G_IMG_H-2)/4, CHANNELS => C3_OUT_CH)
        port map(clk => clk, reset => reset, in_data => s7_data, in_valid => s7_valid, in_ready => s7_ready, out_data => s8_data, out_valid => s8_valid, out_ready => s8_ready);

    -- =========================================================
    -- BRIDGE & CLASSIFIER
    -- =========================================================
    u_flatten: entity work.flatten
        port map(clk => clk, reset => reset, in_data => s8_data, in_valid => s8_valid, in_ready => s8_ready, out_data => s9_data, out_valid => s9_valid, out_ready => s9_ready, flatten_done => flatten_done_open);

    u_dense1: entity work.dense_engine
        generic map(
            IN_SIZE       => D1_IN_SIZE,
            OUT_SIZE      => D1_OUT_SIZE,
            PARALLEL_MACS => 8,
            APPLY_RELU    => true,
            BIAS_FILE     => D1_B_FILE,
            WEIGHT_FILE   => D1_W_FILE
        )
        port map(
            clk => clk, reset => reset,
            in_data => s9_data, in_valid => s9_valid, in_ready => s9_ready,
            out_data => s10_data, out_valid => s10_valid, out_ready => s10_ready
        );

    u_dense2: entity work.dense_engine
        generic map(
            IN_SIZE       => D2_IN_SIZE,
            OUT_SIZE      => D2_OUT_SIZE,
            PARALLEL_MACS => 8,
            APPLY_RELU    => false,
            BIAS_FILE     => D2_B_FILE,
            WEIGHT_FILE   => D2_W_FILE
        )
        port map(
            clk => clk, reset => reset,
            in_data => s10_data, in_valid => s10_valid, in_ready => s10_ready,
            out_data => out_data, out_valid => out_valid, out_ready => out_ready
        );

end architecture;
\end{lstlisting}

\subsection{conv_layer.vhd}
\begin{lstlisting}[language=VHDL]
library IEEE;
use IEEE.STD_LOGIC_1164.ALL;
use IEEE.NUMERIC_STD.ALL;
use work.cnn_config.all;

entity conv_layer is
    generic (
        NUM_IN_CH   : integer := 1; 
        NUM_OUT_CH  : integer := 16;
        IMG_W       : integer := 34;
        IMG_H       : integer := 34;
        OC_PARALLEL : integer := 1;
        WEIGHT_FILE : string  := "conv1_weights.mem";
        BIAS_FILE   : string  := "conv1_biases.mem"
    );
    port (
        clk         : in  std_logic;
        reset       : in  std_logic;
        
        in_data     : in  std_logic_vector(DATA_WIDTH-1 downto 0);
        in_valid    : in  std_logic;
        in_ready    : out std_logic;
        
        out_data    : out std_logic_vector(DATA_WIDTH-1 downto 0);
        out_valid   : out std_logic;
        out_ready   : in  std_logic
    );
end entity;

architecture rtl of conv_layer is
    type state_t is (SHIFT_IN, CALC_MAC, PRE_STREAM, STREAM_OUT);
    signal state : state_t := SHIFT_IN;

    signal col_cnt : integer range 0 to IMG_W - 1 := 0;
    signal row_cnt : integer range 0 to IMG_H - 1 := 0;
    signal ic_cnt  : integer range 0 to NUM_IN_CH - 1 := 0;

    constant BUF_DEPTH : integer := IMG_W * NUM_IN_CH;
    type ram_type is array (0 to BUF_DEPTH - 1) of signed(DATA_WIDTH-1 downto 0);
    signal line_buf_1 : ram_type := (others => (others => '0'));
    signal line_buf_2 : ram_type := (others => (others => '0'));
    signal rd_addr    : integer range 0 to BUF_DEPTH - 1 := 0;

    constant DELAY_LEN : integer := 2 * NUM_IN_CH;
    type delay_line_t is array (0 to DELAY_LEN) of signed(DATA_WIDTH-1 downto 0);
    signal delay_row0 : delay_line_t := (others => (others => '0'));
    signal delay_row1 : delay_line_t := (others => (others => '0'));
    signal delay_row2 : delay_line_t := (others => (others => '0'));

    signal rom_addr : std_logic_vector(31 downto 0) := (others => '0');
    signal rom_dout : std_logic_vector(143 downto 0);
    type weight_array_t is array (0 to 8) of signed(15 downto 0);
    signal real_weights : weight_array_t := (others => (others => '0'));
    
    signal bias_addr : std_logic_vector(31 downto 0) := (others => '0');
    signal bias_dout : std_logic_vector(15 downto 0);
    
    -- NEW: Extended pipeline registers for the Binary Adder Tree
    signal bias_reg_1  : signed(15 downto 0) := (others => '0');
    signal bias_reg_2  : signed(15 downto 0) := (others => '0');
    signal bias_reg_3  : signed(15 downto 0) := (others => '0');
    
    signal sum_A       : signed(47 downto 0) := (others => '0');
    signal sum_B       : signed(47 downto 0) := (others => '0');
    signal sum_C       : signed(47 downto 0) := (others => '0');
    signal mac_sum_reg : signed(47 downto 0) := (others => '0');
    
    type prod_array_t is array (0 to 8) of signed(31 downto 0);
    signal prods : prod_array_t := (others => (others => '0'));

    -- NEW: Expanded pipeline tracker (5 stages now: 0 to 4)
    type pipeline_track_t is array (0 to 4) of integer range -1 to NUM_OUT_CH;
    signal oc_pipe : pipeline_track_t := (others => -1);
    signal oc_req  : integer range 0 to NUM_OUT_CH := 0;
    
    type accum_array_t is array (0 to NUM_OUT_CH - 1) of signed(47 downto 0);
    signal accums : accum_array_t := (others => (others => '0'));
    
    signal out_ch_cnt : integer range 0 to NUM_OUT_CH := 0;

begin

    weight_inst : entity work.weight_rom
        generic map(WORD_WIDTH => 144, DEPTH => NUM_OUT_CH * NUM_IN_CH, INIT_FILE => WEIGHT_FILE)
        port map(clk => clk, addr => rom_addr, dout => rom_dout);

    bias_inst : entity work.weight_rom
        generic map(WORD_WIDTH => 16, DEPTH => NUM_OUT_CH, INIT_FILE => BIAS_FILE)
        port map(clk => clk, addr => bias_addr, dout => bias_dout);

    process(rom_dout)
    begin
        for i in 0 to 8 loop
            real_weights(i) <= signed(rom_dout((i+1)*16-1 downto i*16));
        end loop;
    end process;

    in_ready  <= '1' when state = SHIFT_IN else '0';

    process(clk)
        variable temp_accum   : signed(47 downto 0);
        variable temp_shifted : signed(47 downto 0);
    begin
        if rising_edge(clk) then
            if reset = '1' then
                state <= SHIFT_IN;
                col_cnt <= 0; row_cnt <= 0; ic_cnt <= 0; rd_addr <= 0;
                oc_pipe <= (others => -1); oc_req <= 0; out_ch_cnt <= 0;
                out_valid <= '0';
                out_data <= (others => '0');
            else
                case state is
                    when SHIFT_IN =>
                        out_valid <= '0'; 
                        if in_valid = '1' then
                            line_buf_1(rd_addr) <= signed(in_data);
                            line_buf_2(rd_addr) <= line_buf_1(rd_addr);
                            
                            delay_row0(1 to DELAY_LEN) <= delay_row0(0 to DELAY_LEN - 1); delay_row0(0) <= line_buf_2(rd_addr);
                            delay_row1(1 to DELAY_LEN) <= delay_row1(0 to DELAY_LEN - 1); delay_row1(0) <= line_buf_1(rd_addr);
                            delay_row2(1 to DELAY_LEN) <= delay_row2(0 to DELAY_LEN - 1); delay_row2(0) <= signed(in_data);
                            
                            if rd_addr = BUF_DEPTH - 1 then rd_addr <= 0; else rd_addr <= rd_addr + 1; end if;

                            if row_cnt >= 2 and col_cnt >= 2 then
                                state <= CALC_MAC;
                                oc_req <= 0;
                                oc_pipe <= (others => -1);
                            else
                                if ic_cnt = NUM_IN_CH - 1 then
                                    ic_cnt <= 0;
                                    if col_cnt = IMG_W - 1 then
                                        col_cnt <= 0;
                                        if row_cnt = IMG_H - 1 then row_cnt <= 0; else row_cnt <= row_cnt + 1; end if;
                                    else col_cnt <= col_cnt + 1; end if;
                                else ic_cnt <= ic_cnt + 1; end if;
                            end if;
                        end if;

                    when CALC_MAC =>
                        -- Pipeline Stage 0: Issue Memory Addresses
                        if oc_req < NUM_OUT_CH then
                            rom_addr  <= std_logic_vector(to_unsigned(oc_req * NUM_IN_CH + ic_cnt, 32));
                            bias_addr <= std_logic_vector(to_unsigned(oc_req, 32)); 
                            oc_pipe(0) <= oc_req;
                            oc_req <= oc_req + 1;
                        else
                            oc_pipe(0) <= -1;
                        end if;

                        -- Pipeline Stage 1: Wait for BRAM latency, then Multiply
                        oc_pipe(1) <= oc_pipe(0);
                        if oc_pipe(1) >= 0 then
                            prods(0) <= delay_row0(2 * NUM_IN_CH) * real_weights(0);
                            prods(1) <= delay_row0(NUM_IN_CH)     * real_weights(1);
                            prods(2) <= delay_row0(0)             * real_weights(2);
                            prods(3) <= delay_row1(2 * NUM_IN_CH) * real_weights(3);
                            prods(4) <= delay_row1(NUM_IN_CH)     * real_weights(4);
                            prods(5) <= delay_row1(0)             * real_weights(5);
                            prods(6) <= delay_row2(2 * NUM_IN_CH) * real_weights(6);
                            prods(7) <= delay_row2(NUM_IN_CH)     * real_weights(7);
                            prods(8) <= delay_row2(0)             * real_weights(8);
                            bias_reg_1 <= signed(bias_dout);
                        end if;

                        -- Pipeline Stage 2: Binary Adder Tree (Part 1 - Partial Sums)
                        oc_pipe(2) <= oc_pipe(1);
                        if oc_pipe(2) >= 0 then
                            sum_A <= resize(prods(0),48) + resize(prods(1),48) + resize(prods(2),48) + resize(prods(3),48);
                            sum_B <= resize(prods(4),48) + resize(prods(5),48) + resize(prods(6),48) + resize(prods(7),48);
                            sum_C <= resize(prods(8),48);
                            bias_reg_2 <= bias_reg_1; 
                        end if;

                        -- Pipeline Stage 3: Binary Adder Tree (Part 2 - Final MAC Sum)
                        oc_pipe(3) <= oc_pipe(2);
                        if oc_pipe(3) >= 0 then
                            mac_sum_reg <= sum_A + sum_B + sum_C;
                            bias_reg_3 <= bias_reg_2;
                        end if;

                        -- Pipeline Stage 4: Accumulation
                        oc_pipe(4) <= oc_pipe(3);
                        if oc_pipe(4) >= 0 then
                            if ic_cnt = 0 then
                                accums(oc_pipe(4)) <= mac_sum_reg + shift_left(resize(bias_reg_3, 48), 8); 
                            else
                                accums(oc_pipe(4)) <= accums(oc_pipe(4)) + mac_sum_reg; 
                            end if;
                        end if;

                        -- Flow Control: Triggered on final stage of pipeline (Stage 4)
                        if oc_pipe(4) = NUM_OUT_CH - 1 then
                            if ic_cnt = NUM_IN_CH - 1 then
                                state <= PRE_STREAM;
                                out_ch_cnt <= 0;
                            else
                                state <= SHIFT_IN;
                            end if;
                            
                            if ic_cnt = NUM_IN_CH - 1 then
                                ic_cnt <= 0;
                                if col_cnt = IMG_W - 1 then
                                    col_cnt <= 0;
                                    if row_cnt = IMG_H - 1 then row_cnt <= 0; else row_cnt <= row_cnt + 1; end if;
                                else col_cnt <= col_cnt + 1; end if;
                            else ic_cnt <= ic_cnt + 1; end if;
                        end if;

                    when PRE_STREAM =>
                        temp_accum := accums(out_ch_cnt);
                        temp_shifted := shift_right(temp_accum, 8);
                        if temp_shifted < 0 then
                            out_data <= (others => '0'); 
                        elsif temp_shifted > 32767 then
                            out_data <= std_logic_vector(to_signed(32767, 16)); 
                        else
                            out_data <= std_logic_vector(resize(temp_shifted, 16));
                        end if;
                        out_valid <= '1';
                        state <= STREAM_OUT;

                    when STREAM_OUT =>
                        if out_ready = '1' then
                            if out_ch_cnt = NUM_OUT_CH - 1 then
                                out_valid <= '0';
                                state <= SHIFT_IN;
                            else
                                out_ch_cnt <= out_ch_cnt + 1;
                                temp_accum := accums(out_ch_cnt + 1);
                                temp_shifted := shift_right(temp_accum, 8);
                                if temp_shifted < 0 then
                                    out_data <= (others => '0');
                                elsif temp_shifted > 32767 then
                                    out_data <= std_logic_vector(to_signed(32767, 16));
                                else
                                    out_data <= std_logic_vector(resize(temp_shifted, 16));
                                end if;
                            end if;
                        end if;
                        
                end case;
            end if;
        end if;
    end process;
end architecture;
\end{lstlisting}

\subsection{pool_layer.vhd}
\begin{lstlisting}[language=VHDL]
library IEEE;
use IEEE.STD_LOGIC_1164.ALL;
use IEEE.NUMERIC_STD.ALL;
use work.cnn_config.all;

entity pool_layer is
    generic (
        IMG_W    : integer := 32;
        IMG_H    : integer := 32;
        CHANNELS : integer := 16
    );
    port (
        clk         : in  std_logic;
        reset       : in  std_logic;
        
        in_data     : in  std_logic_vector(DATA_WIDTH-1 downto 0);
        in_valid    : in  std_logic;
        in_ready    : out std_logic;
        
        out_data    : out std_logic_vector(DATA_WIDTH-1 downto 0);
        out_valid   : out std_logic;
        out_ready   : in  std_logic
    );
end entity;

architecture rtl of pool_layer is
    signal col_cnt : integer range 0 to IMG_W - 1 := 0;
    signal row_cnt : integer range 0 to IMG_H - 1 := 0;
    signal ch_cnt  : integer range 0 to CHANNELS - 1 := 0;

    constant BUF_DEPTH : integer := IMG_W * CHANNELS;
    type ram_type is array (0 to BUF_DEPTH - 1) of signed(DATA_WIDTH-1 downto 0);
    signal line_buf : ram_type := (others => (others => '0'));

    type max_array_t is array (0 to CHANNELS - 1) of signed(DATA_WIDTH-1 downto 0);
    signal prev_max_col : max_array_t := (others => (others => '0'));

    -- Pipelining signals (Stage 1 to Stage 2)
    signal current_pixel_reg  : signed(DATA_WIDTH-1 downto 0) := (others => '0');
    signal buffered_pixel_reg : signed(DATA_WIDTH-1 downto 0) := (others => '0');
    signal valid_pipe         : std_logic := '0';
    signal is_output_pipe     : boolean := false;
    signal ch_cnt_pipe        : integer range 0 to CHANNELS - 1 := 0;
    signal col_even_pipe      : boolean := false;

    -- Stage 2 Combinational signals
    signal current_max_col    : signed(DATA_WIDTH-1 downto 0);
    signal is_output_pixel    : boolean;
    signal stall_pipeline     : std_logic;

begin
    -- ========================================================================
    -- Combinational Logic
    -- ========================================================================
    
    -- Identify if the incoming coordinate is a pooling output trigger (Odd row AND Odd col)
    is_output_pixel <= ((row_cnt mod 2) = 1 and (col_cnt mod 2) = 1);

    -- Flow Control: Stall if Stage 2 holds a valid output pixel, but downstream is NOT ready.
    stall_pipeline <= '1' when (valid_pipe = '1' and is_output_pipe and out_ready = '0') else '0';

    -- Input Ready: We can accept new data if the pipeline isn't stalled.
    in_ready <= '1' when (stall_pipeline = '0') else '0';

    -- Stage 2 Combinational Datapath
    current_max_col <= current_pixel_reg when current_pixel_reg > buffered_pixel_reg else buffered_pixel_reg;

    -- Output assignment (Stage 2)
    out_valid <= '1' when (valid_pipe = '1' and is_output_pipe) else '0';
    out_data  <= std_logic_vector(current_max_col) when (current_max_col > prev_max_col(ch_cnt_pipe)) 
                 else std_logic_vector(prev_max_col(ch_cnt_pipe));

    -- ========================================================================
    -- Synchronous Process
    -- ========================================================================
    process(clk)
    begin
        if rising_edge(clk) then
            if reset = '1' then
                col_cnt        <= 0;
                row_cnt        <= 0;
                ch_cnt         <= 0;
                prev_max_col   <= (others => (others => '0'));
                valid_pipe     <= '0';
                is_output_pipe <= false;
                col_even_pipe  <= false;
            else
                -- Proceed only if the pipeline is not stalled
                if stall_pipeline = '0' then
                    
                    -- --------------------------------------------------------
                    -- RETIRE STAGE 2
                    -- --------------------------------------------------------
                    -- Pipeline Left Column Max on even columns using Stage 2 data
                    if valid_pipe = '1' and col_even_pipe then
                        prev_max_col(ch_cnt_pipe) <= current_max_col;
                    end if;

                    -- --------------------------------------------------------
                    -- ADVANCE STAGE 1 -> STAGE 2
                    -- --------------------------------------------------------
                    if in_valid = '1' then
                        -- Latch data and control into pipeline registers
                        buffered_pixel_reg <= line_buf(col_cnt * CHANNELS + ch_cnt);
                        current_pixel_reg  <= signed(in_data);
                        ch_cnt_pipe        <= ch_cnt;
                        is_output_pipe     <= is_output_pixel;
                        col_even_pipe      <= ((col_cnt mod 2) = 0);
                        valid_pipe         <= '1';
                        
                        -- Update Line Buffer on even rows
                        if (row_cnt mod 2) = 0 then
                            line_buf(col_cnt * CHANNELS + ch_cnt) <= signed(in_data);
                        end if;

                        -- Interleaved Counter Logic
                        if ch_cnt = CHANNELS - 1 then
                            ch_cnt <= 0;
                            if col_cnt = IMG_W - 1 then
                                col_cnt <= 0;
                                if row_cnt = IMG_H - 1 then
                                    row_cnt <= 0;
                                else
                                    row_cnt <= row_cnt + 1;
                                end if;
                            else
                                col_cnt <= col_cnt + 1;
                            end if;
                        else
                            ch_cnt <= ch_cnt + 1;
                        end if;
                    else
                        -- No valid input, pipeline bubble
                        valid_pipe <= '0';
                    end if;
                end if;
            end if;
        end if;
    end process;
end architecture;
\end{lstlisting}

\subsection{padder.vhd}
\begin{lstlisting}[language=VHDL]
library IEEE;
use IEEE.STD_LOGIC_1164.ALL;
use IEEE.NUMERIC_STD.ALL;
use work.cnn_config.all;

entity padder is
    generic (
        IMG_W    : integer := 32;
        IMG_H    : integer := 32;
        CHANNELS : integer := 16
    );
    port (
        clk         : in  std_logic;
        reset       : in  std_logic;
        in_data     : in  std_logic_vector(DATA_WIDTH-1 downto 0);
        in_valid    : in  std_logic;
        in_ready    : out std_logic;
        out_data    : out std_logic_vector(DATA_WIDTH-1 downto 0);
        out_valid   : out std_logic;
        out_ready   : in  std_logic
    );
end entity;

architecture rtl of padder is
    type state_type is (IDLE, STREAMING);
    signal state : state_type := IDLE;
    
    -- Counters map the new dimensions (e.g., 18x18 instead of 16x16)
    signal col_cnt : integer range 0 to IMG_W + 1 := 0;
    signal row_cnt : integer range 0 to IMG_H + 1 := 0;
    signal ch_cnt  : integer range 0 to CHANNELS - 1 := 0;
    
    signal is_pad  : std_logic;
begin

    -- Geographical check: Are we in the 1-pixel zero border?
    is_pad <= '1' when (row_cnt = 0 or row_cnt = IMG_H + 1 or col_cnt = 0 or col_cnt = IMG_W + 1) else '0';

    -- Handshake Wiring
    -- We are only ready to accept real data if we are NOT in a padding region AND the downstream is ready
    in_ready  <= out_ready when (state = STREAMING and is_pad = '0') else '0';
    
    -- We are valid if we are streaming padding zeros, OR if we are streaming real data and the input is valid
    out_valid <= '1'       when (state = STREAMING and is_pad = '1') else
                 in_valid  when (state = STREAMING and is_pad = '0') else '0';
                 
    -- Multiplex the output data between zeros and the actual input stream
    out_data  <= (others => '0') when is_pad = '1' else in_data;

    process(clk)
    begin
        if rising_edge(clk) then
            if reset = '1' then
                state   <= IDLE;
                col_cnt <= 0;
                row_cnt <= 0;
                ch_cnt  <= 0;
            else
                case state is
                    when IDLE =>
                        -- Wait for the very first pixel of a new image to arrive before turning on the tap
                        if in_valid = '1' then
                            state <= STREAMING;
                        end if;

                    when STREAMING =>
                        -- Advance counters ONLY when a successful transmission happens
                        if (out_ready = '1' and is_pad = '1') or (out_ready = '1' and in_valid = '1' and is_pad = '0') then
                            
                            if ch_cnt = CHANNELS - 1 then
                                ch_cnt <= 0;
                                
                                if col_cnt = IMG_W + 1 then
                                    col_cnt <= 0;
                                    
                                    if row_cnt = IMG_H + 1 then
                                        row_cnt <= 0;
                                        state   <= IDLE; -- Finished sending the entire padded image
                                    else
                                        row_cnt <= row_cnt + 1;
                                    end if;
                                else
                                    col_cnt <= col_cnt + 1;
                                end if;
                            else
                                ch_cnt <= ch_cnt + 1;
                            end if;
                        end if;
                end case;
            end if;
        end if;
    end process;

end architecture;
\end{lstlisting}

\subsection{flatten.vhd}
\begin{lstlisting}[language=VHDL]
library IEEE;
use IEEE.STD_LOGIC_1164.ALL;
use IEEE.NUMERIC_STD.ALL;
use work.cnn_config.all;

entity flatten is
    port (
        clk          : in  std_logic;
        reset        : in  std_logic;
        
        -- Input Stream (From Pool3)
        in_data      : in  std_logic_vector(DATA_WIDTH-1 downto 0);
        in_valid     : in  std_logic;
        in_ready     : out std_logic;
        
        -- Output Stream (To Dense1)
        out_data     : out std_logic_vector(DATA_WIDTH-1 downto 0);
        out_valid    : out std_logic;
        out_ready    : in  std_logic;
        
        -- Status Flag
        flatten_done : out std_logic
    );
end entity;

architecture rtl of flatten is
    -- Counter sized perfectly for the flatten size defined in your config
    signal count : integer range 0 to FLATTEN_SIZE := 0;
begin
    
    -- 1. Pure Combinatorial Pass-Through
    out_data  <= in_data;
    out_valid <= in_valid;
    in_ready  <= out_ready;

    -- 2. Synchronous Counting Logic
    process(clk)
    begin
        if rising_edge(clk) then
            if reset = '1' then
                count <= 0;
                flatten_done <= '0';
            else
                -- Default state: no pulse
                flatten_done <= '0'; 
                
                -- Only count when data is actively transferring
                if in_valid = '1' and out_ready = '1' then
                    if count = FLATTEN_SIZE - 1 then
                        count <= 0;
                        flatten_done <= '1'; -- Pulse high for 1 cycle when a full map is flattened
                    else
                        count <= count + 1;
                    end if;
                end if;
            end if;
        end if;
    end process;
end architecture;
\end{lstlisting}

\subsection{dense_engine.vhd}
\begin{lstlisting}[language=VHDL]
library IEEE;
use IEEE.STD_LOGIC_1164.ALL;
use IEEE.NUMERIC_STD.ALL;
use std.textio.all;
use ieee.std_logic_textio.all;
use work.cnn_config.all;

entity dense_engine is
    generic (
        IN_SIZE       : integer := 512;
        OUT_SIZE      : integer := 32;
        PARALLEL_MACS : integer := 8;
        APPLY_RELU    : boolean := true;
        BIAS_FILE     : string  := "dense1_biases.mem";
        WEIGHT_FILE   : string  := "dense1_weights.mem"
    );
    port (
        clk         : in  std_logic;
        reset       : in  std_logic;
        
        in_data     : in  std_logic_vector(DATA_WIDTH-1 downto 0);
        in_valid    : in  std_logic;
        in_ready    : out std_logic;
        
        out_data    : out std_logic_vector(DATA_WIDTH-1 downto 0);
        out_valid   : out std_logic;
        out_ready   : in  std_logic
    );
end entity;

architecture rtl of dense_engine is
    type state_type is (IDLE, LOAD_BUFFER, COMPUTE_WAIT, COMPUTE, FLUSH_MAC, DRAIN_OUT);
    signal state : state_type := IDLE;

    -- The Memory Array
    type ram_type is array (0 to IN_SIZE - 1) of signed(DATA_WIDTH-1 downto 0);
    signal input_buf : ram_type := (others => (others => '0'));
    
    attribute ram_style : string;
    attribute ram_style of input_buf : signal is "block";
    
    signal p_in_reg  : signed(DATA_WIDTH-1 downto 0) := (others => '0');
    signal ram_we    : std_logic;

    signal load_idx  : integer range 0 to IN_SIZE := 0;

    constant NUM_PASSES : integer := (OUT_SIZE + PARALLEL_MACS - 1) / PARALLEL_MACS; 
    signal pass_idx     : integer range 0 to NUM_PASSES := 0;
    
    signal fetch_idx    : integer range 0 to IN_SIZE := 0;
    signal compute_idx  : integer range 0 to IN_SIZE := 0;
    
    signal total_outputs_sent : integer range 0 to OUT_SIZE := 0;

    type accum_array is array (0 to PARALLEL_MACS - 1) of signed(47 downto 0);
    signal accums : accum_array := (others => (others => '0'));
    
    signal valid_prod : std_logic := '0';
    signal prods_reg  : accum_array := (others => (others => '0'));

    signal drain_idx : integer range 0 to PARALLEL_MACS := 0;

    type bias_array is array (0 to OUT_SIZE - 1) of signed(15 downto 0);
    
    impure function init_bias(filename : string) return bias_array is
        file f : text;
        variable l : line;
        variable w : std_logic_vector(15 downto 0);
        variable res : bias_array := (others => (others => '0'));
    begin
        file_open(f, filename, read_mode);
        for i in 0 to OUT_SIZE - 1 loop
            if not endfile(f) then
                readline(f, l);
                hread(l, w);
                res(i) := signed(w);
            end if;
        end loop;
        file_close(f);
        return res;
    end function;
    
    signal biases : bias_array := init_bias(BIAS_FILE);

    signal weight_addr : std_logic_vector(31 downto 0);
    signal weight_data : std_logic_vector(PARALLEL_MACS * 16 - 1 downto 0);

begin

    u_weight_rom: entity work.weight_rom
        generic map (
            WORD_WIDTH => PARALLEL_MACS * 16,
            DEPTH      => NUM_PASSES * IN_SIZE,
            INIT_FILE  => WEIGHT_FILE
        )
        port map (
            clk  => clk,
            addr => weight_addr,
            dout => weight_data
        );

    in_ready <= '1' when (state = IDLE or state = LOAD_BUFFER) and load_idx < IN_SIZE else '0';
    
    weight_addr <= std_logic_vector(to_unsigned(pass_idx * IN_SIZE + fetch_idx, 32));

    -- =========================================================
    -- THE FIX: DEDICATED BRAM PROCESS
    -- This perfectly matches the Xilinx synchronous RAM template
    -- =========================================================
    ram_we <= '1' when (state = IDLE or state = LOAD_BUFFER) and in_valid = '1' else '0';

    process(clk)
    begin
        if rising_edge(clk) then
            -- Port A: Write
            if ram_we = '1' then
                input_buf(load_idx) <= signed(in_data);
            end if;
            
            -- Port B: Read (Free-running, perfectly pipeline aligned)
            p_in_reg <= input_buf(fetch_idx);
        end if;
    end process;


    -- =========================================================
    -- MAIN CONTROL FSM
    -- =========================================================
    process(clk)
        variable final_val : signed(47 downto 0);
        variable shifted_val : signed(47 downto 0);
    begin
        if rising_edge(clk) then
            if reset = '1' then
                state <= IDLE;
                load_idx <= 0; pass_idx <= 0; 
                fetch_idx <= 0; compute_idx <= 0;
                drain_idx <= 0;
                valid_prod <= '0';
                total_outputs_sent <= 0;
                accums <= (others => (others => '0'));
                prods_reg <= (others => (others => '0'));
                out_valid <= '0'; out_data <= (others => '0');
            else
                out_valid <= '0'; 

                case state is
                    when IDLE =>
                        load_idx <= 0;
                        total_outputs_sent <= 0;
                        if in_valid = '1' then
                            load_idx <= 1;
                            state <= LOAD_BUFFER;
                        end if;

                    when LOAD_BUFFER =>
                        if in_valid = '1' then
                            if load_idx = IN_SIZE - 1 then
                                state <= COMPUTE_WAIT; 
                                pass_idx <= 0;
                                fetch_idx <= 0;
                                compute_idx <= 0;
                                valid_prod <= '0';
                                accums <= (others => (others => '0'));
                            else
                                load_idx <= load_idx + 1;
                            end if;
                        end if;

                    when COMPUTE_WAIT =>
                        fetch_idx <= 1;
                        compute_idx <= 0;
                        valid_prod <= '0';
                        state <= COMPUTE;

                    when COMPUTE =>
                        for i in 0 to PARALLEL_MACS - 1 loop
                            prods_reg(i) <= resize(p_in_reg * signed(weight_data((i+1)*16-1 downto i*16)), 48); 
                        end loop;
                        valid_prod <= '1';

                        if valid_prod = '1' then
                            for i in 0 to PARALLEL_MACS - 1 loop
                                accums(i) <= accums(i) + prods_reg(i);
                            end loop;
                        end if;

                        if compute_idx = IN_SIZE - 1 then
                            state <= FLUSH_MAC; 
                        else
                            compute_idx <= compute_idx + 1;
                            if fetch_idx < IN_SIZE - 1 then
                                fetch_idx <= fetch_idx + 1;
                            end if;
                        end if;

                    when FLUSH_MAC =>
                        if valid_prod = '1' then
                            for i in 0 to PARALLEL_MACS - 1 loop
                                accums(i) <= accums(i) + prods_reg(i);
                            end loop;
                        end if;
                        valid_prod <= '0';
                        state <= DRAIN_OUT;
                        drain_idx <= 0;

                    when DRAIN_OUT =>
                        if out_ready = '1' then
                            out_valid <= '1';
                            final_val := accums(drain_idx) + shift_left(resize(biases(total_outputs_sent), 48), 8);
                            shifted_val := shift_right(final_val, 8);

                            if APPLY_RELU and shifted_val < 0 then
                                out_data <= (others => '0');
                            elsif shifted_val > 32767 then
                                out_data <= std_logic_vector(to_signed(32767, 16));
                            elsif shifted_val < -32768 then
                                out_data <= std_logic_vector(to_signed(-32768, 16));
                            else
                                out_data <= std_logic_vector(resize(shifted_val, 16));
                            end if;

                            if total_outputs_sent = OUT_SIZE - 1 then
                                state <= IDLE;
                                total_outputs_sent <= 0;
                            elsif drain_idx = PARALLEL_MACS - 1 then
                                if pass_idx < NUM_PASSES - 1 then
                                    state <= COMPUTE_WAIT; 
                                    pass_idx <= pass_idx + 1;
                                    fetch_idx <= 0;
                                    compute_idx <= 0;
                                    accums <= (others => (others => '0'));
                                    drain_idx <= 0;
                                    total_outputs_sent <= total_outputs_sent + 1;
                                else
                                    state <= IDLE;
                                end if;
                            else
                                drain_idx <= drain_idx + 1;
                                total_outputs_sent <= total_outputs_sent + 1;
                            end if;
                        end if;
                end case;
            end if;
        end if;
    end process;
end architecture;
\end{lstlisting}

\subsection{spi_controller.vhd}
\begin{lstlisting}[language=VHDL]
library IEEE;
use IEEE.STD_LOGIC_1164.ALL;

entity spi_controller is
    Generic (
        G_TX_WIDTH : integer := 8;
        G_RX_WIDTH : integer := 8
    );
    Port (
        clk_100mhz : in  STD_LOGIC;
        reset      : in  STD_LOGIC;
        
        -- Asynchronous SPI Pins from Raspberry Pi
        spi_cs_n   : in  STD_LOGIC;
        spi_sclk   : in  STD_LOGIC;
        spi_mosi   : in  STD_LOGIC;
        spi_miso   : out STD_LOGIC;
        
        -- Interface to Top Module
        tx_data    : in  STD_LOGIC_VECTOR(G_TX_WIDTH-1 downto 0);
        tx_load    : in  STD_LOGIC;
        rx_data    : out STD_LOGIC_VECTOR(G_RX_WIDTH-1 downto 0);
        rx_valid   : out STD_LOGIC
    );
end spi_controller;

architecture Behavioral of spi_controller is
    -- Synchronization Registers
    signal cs_sync   : STD_LOGIC_VECTOR(1 downto 0) := "11";
    signal sclk_sync : STD_LOGIC_VECTOR(2 downto 0) := "000"; -- 3 bits for edge detection
    signal mosi_sync : STD_LOGIC_VECTOR(1 downto 0) := "00";

    -- Internal signals
    signal sclk_rise : STD_LOGIC;
    signal sclk_fall : STD_LOGIC;
    signal cs_n_safe : STD_LOGIC;
    signal mosi_safe : STD_LOGIC;

begin

    -- Clock Domain Crossing (CDC) & Edge Detection Process
    process(clk_100mhz)
    begin
        if rising_edge(clk_100mhz) then
            if reset = '1' then
                cs_sync   <= "11";
                sclk_sync <= "000";
                mosi_sync <= "00";
            else
                -- Shift registers for 2-flop synchronization
                cs_sync   <= cs_sync(0) & spi_cs_n;
                sclk_sync <= sclk_sync(1 downto 0) & spi_sclk;
                mosi_sync <= mosi_sync(0) & spi_mosi;
            end if;
        end if;
    end process;

    -- Map safe synchronized signals
    cs_n_safe <= cs_sync(1);
    mosi_safe <= mosi_sync(1);

    -- Edge detection logic using the 3rd flop stage
    sclk_rise <= '1' when (sclk_sync(2 downto 1) = "01") else '0';
    sclk_fall <= '1' when (sclk_sync(2 downto 1) = "10") else '0';

    -- Instantiate Receiver
    Inst_Receiver: entity work.spi_receiver
        generic map ( G_DATA_WIDTH => G_RX_WIDTH )
        port map (
            clk_100mhz => clk_100mhz,
            reset      => reset,
            cs_n_sync  => cs_n_safe,
            sclk_rise  => sclk_rise,
            mosi_sync  => mosi_safe,
            rx_data    => rx_data,
            rx_valid   => rx_valid
        );

    -- Instantiate Transmitter
    Inst_Transmitter: entity work.spi_transmitter
        generic map ( G_DATA_WIDTH => G_TX_WIDTH )
        port map (
            clk_100mhz => clk_100mhz,
            reset      => reset,
            cs_n_sync  => cs_n_safe,
            sclk_fall  => sclk_fall,
            tx_data    => tx_data,
            tx_load    => tx_load,
            miso       => spi_miso
        );

end Behavioral;
\end{lstlisting}

\subsection{spi_receiver.vhd}
\begin{lstlisting}[language=VHDL]
library IEEE;
use IEEE.STD_LOGIC_1164.ALL;
use IEEE.NUMERIC_STD.ALL;

entity spi_receiver is
    Generic ( G_DATA_WIDTH : integer := 32 );
    Port (
        clk_100mhz : in STD_LOGIC;
        reset      : in STD_LOGIC;
        cs_n_sync  : in STD_LOGIC; -- Synchronized Active-Low Chip Select
        sclk_rise  : in STD_LOGIC; -- 1-tick pulse when SCLK goes from 0 to 1
        mosi_sync  : in STD_LOGIC; -- Synchronized MOSI data
        rx_data    : out STD_LOGIC_VECTOR(G_DATA_WIDTH-1 downto 0);
        rx_valid   : out STD_LOGIC -- Pulses HIGH for 1 clock cycle when data is fully received
    );
end spi_receiver;

architecture Behavioral of spi_receiver is
    signal shift_reg : STD_LOGIC_VECTOR(G_DATA_WIDTH-1 downto 0) := (others => '0');
    signal bit_count : integer range 0 to G_DATA_WIDTH := 0;
begin
    process(clk_100mhz)
    begin
        if rising_edge(clk_100mhz) then
            if reset = '1' or cs_n_sync = '1' then
                -- CS going High acts as a hard reset for the SPI transaction
                bit_count <= 0;
                rx_valid <= '0';
            else
                rx_valid <= '0'; -- Default state
                
                if sclk_rise = '1' then
                    -- Shift data to the left (MSB first)
                    shift_reg <= shift_reg(G_DATA_WIDTH-2 downto 0) & mosi_sync;
                    bit_count <= bit_count + 1;
                    
                    -- Check if we received the full word
                    if bit_count = G_DATA_WIDTH - 1 then
                        rx_data <= shift_reg(G_DATA_WIDTH-2 downto 0) & mosi_sync;
                        rx_valid <= '1';
                        bit_count <= 0; -- Reset for the next word
                    end if;
                end if;
            end if;
        end if;
    end process;
end Behavioral;
\end{lstlisting}

\subsection{spi_transmitter.vhd}
\begin{lstlisting}[language=VHDL]
library IEEE;
use IEEE.STD_LOGIC_1164.ALL;
use IEEE.NUMERIC_STD.ALL;

entity spi_transmitter is
    Generic ( G_DATA_WIDTH : integer := 32 );
    Port (
        clk_100mhz : in STD_LOGIC;
        reset      : in STD_LOGIC;
        cs_n_sync  : in STD_LOGIC;
        sclk_fall  : in STD_LOGIC; -- 1-tick pulse when SCLK goes from 1 to 0
        tx_data    : in STD_LOGIC_VECTOR(G_DATA_WIDTH-1 downto 0);
        tx_load    : in STD_LOGIC; -- Tells the module to latch new data to send
        miso       : out STD_LOGIC
    );
end spi_transmitter;

architecture Behavioral of spi_transmitter is
    signal shift_reg : STD_LOGIC_VECTOR(G_DATA_WIDTH-1 downto 0) := (others => '0');
    signal miso_internal : STD_LOGIC := '0';
begin
    process(clk_100mhz)
    begin
        if rising_edge(clk_100mhz) then
            if reset = '1' then
                shift_reg <= (others => '0');
            elsif tx_load = '1' then
                -- Load the new data from the Top Module when ready
                shift_reg <= tx_data;
            elsif cs_n_sync = '0' then
                -- When CS is low, shift data out on the falling edge
                if sclk_fall = '1' then
                    shift_reg <= shift_reg(G_DATA_WIDTH-2 downto 0) & '0';
                end if;
            end if;
        end if;
    end process;

    -- Push the Most Significant Bit (MSB) to the wire
    miso_internal <= shift_reg(G_DATA_WIDTH-1);
    
    -- Tristate MISO when CS is high (standard SPI practice)
    miso <= miso_internal when cs_n_sync = '0' else 'Z';

end Behavioral;
\end{lstlisting}

\subsection{sync_fifo.vhd}
\begin{lstlisting}[language=VHDL]
library IEEE;
use IEEE.STD_LOGIC_1164.ALL;
use IEEE.NUMERIC_STD.ALL;

entity sync_fifo is
    generic (
        DATA_WIDTH : integer := 8;
        ADDR_WIDTH : integer := 4 -- 16-byte buffer is plenty for SPI speeds
    );
    port (
        clk      : in  std_logic;
        reset    : in  std_logic;
        wr_en    : in  std_logic;
        wr_data  : in  std_logic_vector(DATA_WIDTH-1 downto 0);
        rd_en    : in  std_logic;
        rd_data  : out std_logic_vector(DATA_WIDTH-1 downto 0);
        empty    : out std_logic;
        full     : out std_logic
    );
end entity;

architecture rtl of sync_fifo is
    type ram_type is array (0 to (2**ADDR_WIDTH)-1) of std_logic_vector(DATA_WIDTH-1 downto 0);
    signal mem : ram_type := (others => (others => '0'));
    
    signal wr_ptr, rd_ptr : unsigned(ADDR_WIDTH-1 downto 0) := (others => '0');
    signal count : integer range 0 to 2**ADDR_WIDTH := 0;
begin
    full  <= '1' when count = 2**ADDR_WIDTH else '0';
    empty <= '1' when count = 0 else '0';
    
    rd_data <= mem(to_integer(rd_ptr));

    process(clk)
    begin
        if rising_edge(clk) then
            if reset = '1' then
                wr_ptr <= (others => '0');
                rd_ptr <= (others => '0');
                count  <= 0;
            else
                if wr_en = '1' and count < 2**ADDR_WIDTH then
                    mem(to_integer(wr_ptr)) <= wr_data;
                    wr_ptr <= wr_ptr + 1;
                end if;
                
                if rd_en = '1' and count > 0 then
                    rd_ptr <= rd_ptr + 1;
                end if;
                
                if (wr_en = '1' and count < 2**ADDR_WIDTH) and not (rd_en = '1' and count > 0) then
                    count <= count + 1;
                elsif (rd_en = '1' and count > 0) and not (wr_en = '1' and count < 2**ADDR_WIDTH) then
                    count <= count - 1;
                end if;
            end if;
        end if;
    end process;
end architecture;
\end{lstlisting}

\subsection{weight_rom.vhd}
\begin{lstlisting}[language=VHDL]
library IEEE;
use IEEE.STD_LOGIC_1164.ALL;
use IEEE.NUMERIC_STD.ALL;
use std.textio.all;
use ieee.std_logic_textio.all;

entity weight_rom is
    generic (
        WORD_WIDTH : integer := 144; -- 9 weights * 16 bits
        DEPTH      : integer := 16;  -- Number of filters/kernels
        INIT_FILE  : string  := "conv1_weights.mem"
    );
    port (
        clk  : in  std_logic;
        addr : in  std_logic_vector(31 downto 0); 
        dout : out std_logic_vector(WORD_WIDTH-1 downto 0)
    );
end entity;

architecture rtl of weight_rom is
    type rom_type is array (0 to DEPTH-1) of std_logic_vector(WORD_WIDTH-1 downto 0);

    -- This function reads 9 lines from your script's output to fill 1 hardware word
    impure function init_rom_from_16bit_file(file_name : string) return rom_type is
        file romfile       : text;
        variable romline   : line;
        variable temp_mem  : rom_type := (others => (others => '0'));
        variable temp_16bit : std_logic_vector(15 downto 0);
        constant WORDS_PER_LINE : integer := WORD_WIDTH / 16;
    begin
        file_open(romfile, file_name, read_mode);
        for i in 0 to DEPTH-1 loop
            for j in 0 to WORDS_PER_LINE-1 loop
                if not endfile(romfile) then
                    readline(romfile, romline);
                    hread(romline, temp_16bit); -- Reads hex like "0100"
                    -- Packs it into the wide 144-bit kernel
                    temp_mem(i)((j+1)*16-1 downto j*16) := temp_16bit;
                end if;
            end loop;
        end loop;
        file_close(romfile);
        return temp_mem;
    end function;

    signal rom_data : rom_type := init_rom_from_16bit_file(INIT_FILE);
    
    -- Forces Vivado to use Block RAM as per Master Plan
    attribute ram_style : string;
    attribute ram_style of rom_data : signal is "block";
begin
    process(clk) begin
        if rising_edge(clk) then
            dout <= rom_data(to_integer(unsigned(addr)));
        end if;
    end process;
end architecture;
\end{lstlisting}

\subsection{eth_cnn_bridge.vhd}
\begin{lstlisting}[language=VHDL]
library IEEE;
use IEEE.STD_LOGIC_1164.ALL;
use IEEE.NUMERIC_STD.ALL;
use work.cnn_config.all; -- Grabs NUM_CLASSES

entity eth_cnn_bridge is
    Generic (
        RESPONSE_LEN : integer := NUM_CLASSES * 2 -- Automatically scales to 6 bytes
    );
    Port (
        clk           : in STD_LOGIC;
        reset_n       : in STD_LOGIC;
        
        s_rx_tdata    : in STD_LOGIC_VECTOR(7 downto 0);
        s_rx_tvalid   : in STD_LOGIC;
        s_rx_tlast    : in STD_LOGIC;
        s_rx_tready   : out STD_LOGIC;
        
        m_tx_tdata    : out STD_LOGIC_VECTOR(7 downto 0);
        m_tx_tvalid   : out STD_LOGIC;
        m_tx_tlast    : out STD_LOGIC;
        m_tx_tuser    : out STD_LOGIC_VECTOR(11 downto 0);
        m_tx_tready   : in  STD_LOGIC
    );
end eth_cnn_bridge;

architecture Behavioral of eth_cnn_bridge is

    -- Reset signals
    signal internal_fifo_reset : std_logic := '0';
    signal fifo_reset_combined : std_logic;

    signal rx_fifo_wr_en, rx_fifo_rd_en : std_logic;
    signal rx_fifo_full, rx_fifo_empty  : std_logic;
    signal rx_fifo_rd_data : std_logic_vector(7 downto 0);

    signal cnn_in_data   : STD_LOGIC_VECTOR(15 downto 0);
    signal cnn_in_valid, cnn_in_ready  : STD_LOGIC;
    signal cnn_out_data  : STD_LOGIC_VECTOR(15 downto 0);
    signal cnn_out_valid, cnn_out_ready : STD_LOGIC;

    type tx_state_type is (IDLE, SEND_RESP);
    signal tx_state : tx_state_type := IDLE;
    signal tx_byte_count : integer range 0 to RESPONSE_LEN := 0;

    signal class_idx : integer range 0 to NUM_CLASSES := 0;
    type byte_array is array (0 to (NUM_CLASSES * 2) - 1) of std_logic_vector(7 downto 0);
    signal tx_buffer : byte_array := (others => x"00");

begin
    -- The FIFO resets if a global hardware reset occurs OR an automatic flush pulse triggers
    fifo_reset_combined <= (not reset_n) or internal_fifo_reset;

    -- RX FIFO
    u_rx_fifo: entity work.sync_fifo
        generic map ( DATA_WIDTH => 8, ADDR_WIDTH => 11 )
        port map ( 
            clk     => clk, 
            reset   => fifo_reset_combined, -- Modified port map connection
            wr_en   => rx_fifo_wr_en, 
            wr_data => s_rx_tdata, 
            rd_en   => rx_fifo_rd_en, 
            rd_data => rx_fifo_rd_data, 
            empty   => rx_fifo_empty, 
            full    => rx_fifo_full 
        );

    rx_fifo_wr_en <= s_rx_tvalid and (not rx_fifo_full);
    s_rx_tready   <= not rx_fifo_full;
    rx_fifo_rd_en <= cnn_in_ready and (not rx_fifo_empty);
    cnn_in_valid  <= not rx_fifo_empty;
    
    cnn_in_data   <= x"00" & rx_fifo_rd_data;

    -- CNN CORE
    u_cnn_core: entity work.cnn_datapath_top
        port map( 
            clk       => clk, 
            reset     => fifo_reset_combined, -- Core resets synchronously with the FIFO
            in_data   => cnn_in_data, 
            in_valid  => cnn_in_valid, 
            in_ready  => cnn_in_ready, 
            out_data  => cnn_out_data, 
            out_valid => cnn_out_valid, 
            out_ready => cnn_out_ready 
        );
        
    cnn_out_ready <= '1'; 

    m_tx_tuser <= std_logic_vector(to_unsigned(RESPONSE_LEN, 12));

    process(clk)
    begin
        if rising_edge(clk) then
            if reset_n = '0' then
                tx_state            <= IDLE;
                tx_byte_count       <= 0;
                class_idx           <= 0;
                internal_fifo_reset <= '0';
            else
                case tx_state is
                    when IDLE =>
                        internal_fifo_reset <= '0'; -- Deassert flush pulse immediately on entering IDLE
                        
                        if cnn_out_valid = '1' then
                            tx_buffer(class_idx * 2)     <= cnn_out_data(15 downto 8);
                            tx_buffer(class_idx * 2 + 1) <= cnn_out_data(7 downto 0);
                            
                            if class_idx = NUM_CLASSES - 1 then
                                class_idx <= 0;
                                tx_state  <= SEND_RESP;
                            else
                                class_idx <= class_idx + 1;
                            end if;
                        end if;

                    when SEND_RESP =>
                        if m_tx_tready = '1' then
                            if tx_byte_count = RESPONSE_LEN - 1 then
                                tx_state            <= IDLE;
                                tx_byte_count       <= 0;
                                internal_fifo_reset <= '1'; -- Pulse active high to flush leftover bytes
                            else
                                tx_byte_count       <= tx_byte_count + 1;
                            end if;
                        end if;
                end case;
            end if;
        end if;
    end process;

    m_tx_tdata  <= tx_buffer(tx_byte_count);
    m_tx_tvalid <= '1' when tx_state = SEND_RESP else '0';
    m_tx_tlast  <= '1' when (tx_state = SEND_RESP and tx_byte_count = RESPONSE_LEN - 1) else '0';

end Behavioral;
\end{lstlisting}

\subsection{tb_system.vhd}
\begin{lstlisting}[language=VHDL]
library IEEE;
use IEEE.STD_LOGIC_1164.ALL;
use IEEE.NUMERIC_STD.ALL;
use std.textio.all;
use ieee.std_logic_textio.all;

entity tb_system is
end entity;

architecture sim of tb_system is

    -- =========================================================
    -- TEST CONFIGURATION
    -- =========================================================
    -- IMPORTANT: Set this to the absolute path where generate_testbench_data.py saves files.
    -- Use forward slashes (/) for VHDL paths, even on Windows.
    constant DATA_DIR   : string := "../../../../testbench_data/"; 
    
    -- Change this prefix to test different images (e.g., "cat", "dog", "human")
    constant TEST_IMAGE : string := "cat"; 

    -- =========================================================
    -- CLOCKS & TIMEOUTS
    -- =========================================================
    constant SYS_CLK_PERIOD : time := 10 ns;   -- 100 MHz FPGA Clock
    constant SPI_CLK_HALF   : time := 50 ns;   -- 10 MHz SPI Clock
    constant TIMEOUT        : time := 50 ms;   -- Failsafe timeout

    -- UUT Signals
    signal clk_100mhz : std_logic := '0';
    signal reset_n    : std_logic := '0';
    signal spi_cs_n   : std_logic := '1';
    signal spi_sclk   : std_logic := '0';
    signal spi_mosi   : std_logic := '0';
    signal spi_miso   : std_logic;
    signal data_ready : std_logic;

    -- Simulation Tracking
    signal sim_done : boolean := false;

    -- 3-Class Arrays for Verification
    type logit_array is array (0 to 2) of std_logic_vector(15 downto 0);
    signal expected_logits : logit_array := (others => (others => '0'));
    signal received_logits : logit_array := (others => (others => '0'));
    -- Add this signal right below your other signals
    signal led : std_logic_vector(15 downto 0);

begin

    -- 100 MHz System Clock Generation
    clk_100mhz <= not clk_100mhz after SYS_CLK_PERIOD/2 when not sim_done else '0';

    -- ... Inside your architecture ...
    uut: entity work.system_top
        port map (
            clk_100mhz => clk_100mhz,
            reset_n    => reset_n,
            spi_cs_n   => spi_cs_n,
            spi_sclk   => spi_sclk,
            spi_mosi   => spi_mosi,
            spi_miso   => spi_miso,
            data_ready => data_ready,
            led        => led  -- Mapped the new port!
        );

    -- =========================================================
    -- MAIN STIMULUS (THE RASPBERRY PI ROLEPLAY)
    -- =========================================================
    pi_sim_proc: process
        file f_expected  : text;
        file f_in        : text;
        variable l       : line;
        variable hex_val : std_logic_vector(15 downto 0);
        variable rx_word : std_logic_vector(15 downto 0);
        variable errors  : integer := 0;
        
        -- Inline Procedure to act as the SPI Master (Mode 0: CPOL=0, CPHA=0)
        procedure spi_transfer_16(
            tx_data : in  std_logic_vector(15 downto 0);
            rx_data : out std_logic_vector(15 downto 0)
        ) is
        begin
            for i in 15 downto 0 loop
                -- 1. Setup data on MOSI
                spi_mosi <= tx_data(i);
                wait for SPI_CLK_HALF;
                
                -- 2. Rising Edge (Sample MISO)
                spi_sclk <= '1';
                rx_data(i) := spi_miso;
                wait for SPI_CLK_HALF;
                
                -- 3. Falling Edge
                spi_sclk <= '0';
            end loop;
        end procedure;

    begin
        report "==================================================";
        report "--- STARTING SYSTEM LEVEL SPI TESTBENCH ---";
        report "Testing Image Prefix: " & TEST_IMAGE;
        report "==================================================";

        -- 1. Load Expected Golden Logits
        report "Loading expected logits from Golden Model...";
        file_open(f_expected, DATA_DIR & TEST_IMAGE & "_11_logits.hex", read_mode);
        for i in 0 to 2 loop
            readline(f_expected, l);
            hread(l, hex_val);
            expected_logits(i) <= hex_val;
        end loop;
        file_close(f_expected);

        -- 2. Hardware Reset
        reset_n  <= '0';
        spi_cs_n <= '1';
        spi_sclk <= '0';
        spi_mosi <= '0';
        wait for 200 ns;
        reset_n  <= '1';
        wait for 200 ns;

        -- 3. Start SPI Transaction
        report "Initiating SPI Transaction (Lowering CS_N)...";
        spi_cs_n <= '0';
        wait for SPI_CLK_HALF * 2;

        -- 4. Transmit the Image (1156 Pixels)
        file_open(f_in, DATA_DIR & TEST_IMAGE & "_0_input.hex", read_mode);
        report "Streaming 1156 Image Words over SPI...";
        
        while not endfile(f_in) loop
            readline(f_in, l);
            hread(l, hex_val);
            spi_transfer_16(hex_val, rx_word); -- Ignore rx_word during image send
        end loop;
        file_close(f_in);
        
        report "Transmission Complete. Waiting for FPGA DATA_READY...";

        -- 5. Wait for CNN Inference (Timeout protection)
        wait until data_ready = '1' for TIMEOUT;
        if data_ready = '0' then
            report "!!! FATAL ERROR: Timeout waiting for DATA_READY !!!" severity failure;
        end if;
        report "DATA_READY Interrupt Received! Reading 3 Class Logits...";

        wait for SPI_CLK_HALF * 4;

        -- 6. Read the 3 Output Logits (Cat, Dog, Human)
        for i in 0 to 2 loop
            spi_transfer_16(x"0000", rx_word);
            received_logits(i) <= rx_word;
        end loop;

        -- 7. End Transaction
        wait for SPI_CLK_HALF * 2;
        spi_cs_n <= '1';
        report "SPI Transaction Closed.";

        -- 8. Bit-Exact Verification
        report "==================================================";
        report "--- VERIFICATION RESULTS ---";
        for i in 0 to 2 loop
            if received_logits(i) = expected_logits(i) then
                report "Logit " & integer'image(i) & " MATCH! Expected/Received: " & to_hex_string(expected_logits(i));
            else
                report "Logit " & integer'image(i) & " MISMATCH! " & 
                       " Expected: " & to_hex_string(expected_logits(i)) & 
                       " Received: " & to_hex_string(received_logits(i)) severity error;
                errors := errors + 1;
            end if;
        end loop;

        report "==================================================";
        if errors = 0 then
            report " >>> SUCCESS: ALL LOGITS MATCH GOLDEN MODEL EXACTLY! <<< ";
        else
            report " >>> FAILURE: " & integer'image(errors) & " LOGIT(S) MISMATCHED <<< " severity error;
        end if;
        report "==================================================";

        sim_done <= true;
        wait;
    end process;

end architecture;
\end{lstlisting}

\subsection{nexys4_top.xdc}
\begin{lstlisting}[language=Tcl]
## ==============================================================================
## FPGA-Accelerated CNN Inference - Master Constraints File
## Target Board: Digilent Nexys 4 / Nexys A7 (Artix-7 XC7A100T)
## ==============================================================================

## ------------------------------------------------------------------------------
## 1. System Clock (100 MHz Oscillator)
## ------------------------------------------------------------------------------
# Mapped to physical pin E3 (100MHz onboard oscillator)
set_property -dict { PACKAGE_PIN E3    IOSTANDARD LVCMOS33 } [get_ports { clk_100mhz }];

# Define the clock for the timing analyzer (10.00 ns period = 100 MHz)
create_clock -add -name sys_clk_pin -period 10.00 -waveform {0 5} [get_ports { clk_100mhz }];


## ------------------------------------------------------------------------------
## 2. System Reset
## ------------------------------------------------------------------------------
# Mapped to the red CPU_RESET button (Active Low)
set_property -dict { PACKAGE_PIN C12   IOSTANDARD LVCMOS33 } [get_ports { reset_n }];


## ------------------------------------------------------------------------------
## 3. SPI Communication Interface (PMOD JA)
## ------------------------------------------------------------------------------
# JA[1] - MOSI (Pi -> FPGA)
set_property -dict { PACKAGE_PIN C17   IOSTANDARD LVCMOS33 } [get_ports { spi_mosi }];

# JA[2] - MISO (FPGA -> Pi)
set_property -dict { PACKAGE_PIN D18   IOSTANDARD LVCMOS33 } [get_ports { spi_miso }];

# JA[3] - SCLK (Pi -> FPGA)
set_property -dict { PACKAGE_PIN E18   IOSTANDARD LVCMOS33 } [get_ports { spi_sclk }];

# JA[4] - CS_N / CE0 (Pi -> FPGA, Active Low)
set_property -dict { PACKAGE_PIN G17   IOSTANDARD LVCMOS33 } [get_ports { spi_cs_n }];


## ------------------------------------------------------------------------------
## 4. Hardware Interrupt (PMOD JB)
## ------------------------------------------------------------------------------
# JB[1] - DATA_READY (FPGA -> Pi)
set_property -dict { PACKAGE_PIN D14   IOSTANDARD LVCMOS33 } [get_ports { data_ready }];


## ------------------------------------------------------------------------------
## 5. Asynchronous Input Constraints
## ------------------------------------------------------------------------------
# Since spi_mosi, spi_sclk, spi_cs_n, and reset_n are external signals asynchronous 
# to the 100MHz clock, we tell the timing analyzer to ignore the setup/hold times 
# from the I/O pads to the first synchronizer flip-flops in spi_controller.vhd.
set_false_path -from [get_ports { spi_mosi spi_sclk spi_cs_n reset_n }]


## ------------------------------------------------------------------------------
## 6. Hardware Debugger LEDs
## ------------------------------------------------------------------------------
# LEDs 15-12: FSM State | LED 11: RX FIFO Empty | LEDs 10-0: Pixel Count
set_property -dict { PACKAGE_PIN H17   IOSTANDARD LVCMOS33 } [get_ports { led[0] }];
set_property -dict { PACKAGE_PIN K15   IOSTANDARD LVCMOS33 } [get_ports { led[1] }];
set_property -dict { PACKAGE_PIN J13   IOSTANDARD LVCMOS33 } [get_ports { led[2] }];
set_property -dict { PACKAGE_PIN N14   IOSTANDARD LVCMOS33 } [get_ports { led[3] }];
set_property -dict { PACKAGE_PIN R18   IOSTANDARD LVCMOS33 } [get_ports { led[4] }];
set_property -dict { PACKAGE_PIN V17   IOSTANDARD LVCMOS33 } [get_ports { led[5] }];
set_property -dict { PACKAGE_PIN U17   IOSTANDARD LVCMOS33 } [get_ports { led[6] }];
set_property -dict { PACKAGE_PIN U16   IOSTANDARD LVCMOS33 } [get_ports { led[7] }];
set_property -dict { PACKAGE_PIN V16   IOSTANDARD LVCMOS33 } [get_ports { led[8] }];
set_property -dict { PACKAGE_PIN T15   IOSTANDARD LVCMOS33 } [get_ports { led[9] }];
set_property -dict { PACKAGE_PIN U14   IOSTANDARD LVCMOS33 } [get_ports { led[10] }];
set_property -dict { PACKAGE_PIN T16   IOSTANDARD LVCMOS33 } [get_ports { led[11] }];
set_property -dict { PACKAGE_PIN V15   IOSTANDARD LVCMOS33 } [get_ports { led[12] }];
set_property -dict { PACKAGE_PIN V14   IOSTANDARD LVCMOS33 } [get_ports { led[13] }];
set_property -dict { PACKAGE_PIN V12   IOSTANDARD LVCMOS33 } [get_ports { led[14] }];
set_property -dict { PACKAGE_PIN V11   IOSTANDARD LVCMOS33 } [get_ports { led[15] }];
\end{lstlisting}


\section{Python Scripts}

\subsection{train_cnn.py}
\begin{lstlisting}[language=Python]
import tensorflow as tf
from tensorflow.keras import layers, models

# ─────────────────────────────────────────────
#  3-CLASS MODEL TRAINER: Cat / Dog / Human
# ─────────────────────────────────────────────

print("Loading images from Doglets/...")
dataset = tf.keras.utils.image_dataset_from_directory(
    "Doglets",
    color_mode="grayscale",
    image_size=(32, 32),
    batch_size=32,
    shuffle=True,
    seed=42
)

class_names = dataset.class_names
print(f"Class mapping: {list(enumerate(class_names))}")

# Split off 20% for validation
total_batches = tf.data.experimental.cardinality(dataset).numpy()
val_size   = int(total_batches * 0.2)
train_size = total_batches - val_size

train_ds = dataset.take(train_size).cache().prefetch(tf.data.AUTOTUNE)
val_ds   = dataset.skip(train_size).cache().prefetch(tf.data.AUTOTUNE)

# ─────────────────────────────────────────────
#  Build the brain (Updated with Normalization)
# ─────────────────────────────────────────────
model = models.Sequential([
    layers.Input(shape=(32, 32, 1)),
    layers.Rescaling(1./255), # Normalizes 0-255 pixels to 0.0-1.0
    
    layers.Conv2D(16, (3, 3), activation='relu', padding='same'),
    layers.MaxPooling2D((2, 2)),
    
    layers.Conv2D(32, (3, 3), activation='relu', padding='same'),
    layers.MaxPooling2D((2, 2)),
    
    layers.Conv2D(32, (3, 3), activation='relu', padding='same'),
    layers.MaxPooling2D((2, 2)),
    
    layers.Flatten(),
    layers.Dense(32, activation='relu'),
    layers.Dense(3, activation='softmax')
])

model.compile(
    optimizer='adam',
    loss='sparse_categorical_crossentropy',
    metrics=['accuracy']
)

# ─────────────────────────────────────────────
#  Train (Increased to 30 Epochs for convergence)
# ─────────────────────────────────────────────
print("\nStarting training (30 epochs)...")
model.fit(train_ds, validation_data=val_ds, epochs=50)

model.save("tiny_ai_brain.h5")
print("\nTraining complete! Saved as tiny_ai_brain.h5")
\end{lstlisting}

\subsection{export_weights.py}
\begin{lstlisting}[language=Python]
import os
import numpy as np
import tensorflow as tf

# ─────────────────────────────────────────────
# 1. Configuration & Helper Functions
# ─────────────────────────────────────────────
MODEL_PATH = "tiny_ai_brain.h5"
OUTPUT_DIR = "export_q88"

def quantize_q8_8(x_float):
    """Converts a float array to Q8.8 fixed-point 16-bit integers."""
    x_int = np.round(x_float * 256.0)
    x_clipped = np.clip(x_int, -32768, 32767)
    return x_clipped.astype(np.int16)

def to_hex_16bit(val):
    """Converts an int16 to a 4-char hex string (2's complement)."""
    return f"{int(val) & 0xFFFF:04X}"

def write_mem_file(filepath, data_1d):
    """Writes a .mem file for Vivado simulation ($readmemh)."""
    with open(filepath, 'w') as f:
        for val in data_1d:
            f.write(to_hex_16bit(val) + "\n")

def write_coe_file(filepath, data_1d):
    """Writes a .coe file for Xilinx Block Memory Generator IP."""
    with open(filepath, 'w') as f:
        f.write("memory_initialization_radix=16;\n")
        f.write("memory_initialization_vector=\n")
        # Write all but the last with a comma
        for val in data_1d[:-1]:
            f.write(to_hex_16bit(val) + ",\n")
        # Write the last one with a semicolon
        f.write(to_hex_16bit(data_1d[-1]) + ";\n")

# ─────────────────────────────────────────────
# 2. Setup & Model Loading
# ─────────────────────────────────────────────
if not os.path.exists(OUTPUT_DIR):
    os.makedirs(OUTPUT_DIR)

print(f"Loading model: {MODEL_PATH}...")
model = tf.keras.models.load_model(MODEL_PATH)

conv_counter = 1
dense_counter = 1

# ─────────────────────────────────────────────
# 3. Layer Iteration & Extraction
# ─────────────────────────────────────────────
for layer in model.layers:
    weights_list = layer.get_weights()
    if len(weights_list) == 0:
        continue # Skip Flatten, MaxPooling, etc.

    w_float = weights_list[0]
    b_float = weights_list[1]
    
    layer_type = layer.__class__.__name__
    
    if layer_type == "Conv2D":
        name = f"conv{conv_counter}"
        conv_counter += 1
        
        # Keras format: (H, W, IC, OC) -> FPGA format: (OC, IC, H, W)
        w_float_transposed = np.transpose(w_float, (3, 2, 0, 1))
        
        # Quantize
        w_q88 = quantize_q8_8(w_float_transposed)
        b_q88 = quantize_q8_8(b_float)
        
    elif layer_type == "Dense":
        name = f"dense{dense_counter}"
        dense_counter += 1
        
        # Initial quantization
        w_q88_base = quantize_q8_8(w_float)
        b_q88_base = quantize_q8_8(b_float)
        
        # Hardware formatting: (Passes, Inputs, Parallel_MACs)
        PARALLEL_MACS = 8
        in_size, out_size = w_q88_base.shape
        num_passes = int(np.ceil(out_size / PARALLEL_MACS))
        pad_cols = (num_passes * PARALLEL_MACS) - out_size
        
        # Pad outputs with zeros if not a perfect multiple of 8 (e.g., Dense2)
        if pad_cols > 0:
            w_q88_padded = np.pad(w_q88_base, ((0,0), (0, pad_cols)), mode='constant')
            b_q88_padded = np.pad(b_q88_base, (0, pad_cols), mode='constant')
        else:
            w_q88_padded = w_q88_base
            b_q88_padded = b_q88_base
            
        # Reshape to (Inputs, Passes, MACs), then transpose to (Passes, Inputs, MACs)
        w_reshaped = w_q88_padded.reshape((in_size, num_passes, PARALLEL_MACS))
        w_hw_format = w_reshaped.transpose((1, 0, 2))
        
        # Assign to w_q88 / b_q88 so the rest of the script exports them correctly
        w_q88 = w_hw_format
        b_q88 = b_q88_padded
        
    else:
        print(f"Skipping unsupported layer with weights: {layer_type}")
        continue
        
    # Flatten arrays for 1D file writing
    w_flat = w_q88.flatten()
    b_flat = b_q88.flatten()

    print(f"Exporting {name}: Weights shape {w_q88.shape} ({len(w_flat)} words), Biases shape {b_q88.shape}")

    # 1. Save .npy for the Golden Python Model
    np.save(os.path.join(OUTPUT_DIR, f"{name}_weights.npy"), w_q88)
    np.save(os.path.join(OUTPUT_DIR, f"{name}_biases.npy"), b_q88)
    
    # 2. Save .mem files (for VHDL simulation)
    write_mem_file(os.path.join(OUTPUT_DIR, f"{name}_weights.mem"), w_flat)
    write_mem_file(os.path.join(OUTPUT_DIR, f"{name}_biases.mem"), b_flat)
    
    # 3. Save .coe files (for Vivado BRAM Synthesis)
    write_coe_file(os.path.join(OUTPUT_DIR, f"{name}_weights.coe"), w_flat)
    write_coe_file(os.path.join(OUTPUT_DIR, f"{name}_biases.coe"), b_flat)

print(f"\nSuccess! All weights exported to the '{OUTPUT_DIR}/' directory.")
\end{lstlisting}

\subsection{golden_model.py}
\begin{lstlisting}[language=Python]
import numpy as np
import os
from PIL import Image

# golden_model.py
# Master Plan Phase 1: Integer-Only Reference Model
# Updated: Input pixels are now properly scaled to Q8.8 to match spi.py hardware inputs!

# ─────────────────────────────────────────────
#  Phase 1.1: Core Arithmetic Engine (Bit-Exact Q8.8)
# ─────────────────────────────────────────────
def dsp48_mac(inputs, weights, bias, apply_relu=True):
    """Simulates the FPGA DSP48E1 48-bit MAC engine."""
    accumulator = 0 
    
    for i, w in zip(inputs, weights):
        accumulator += int(i) * int(w)
        
    accumulator += (int(bias) << 8)
    result = accumulator >> 8
    
    if result > 32767:
        result = 32767
    elif result < -32768:
        result = -32768
        
    if apply_relu and result < 0:
        result = 0
        
    return result

# ─────────────────────────────────────────────
#  Phase 1.2: Spatial & Dense Operations
# ─────────────────────────────────────────────
def conv2d_valid(input_map, weights, biases, apply_relu=True):
    in_h = len(input_map)
    in_w = len(input_map[0])
    in_c = len(input_map[0][0])
    out_c = len(weights)
    
    out_h = in_h - 2
    out_w = in_w - 2
    
    out_map = [[[0 for _ in range(out_c)] for _ in range(out_w)] for _ in range(out_h)]
    
    for y in range(out_h):
        for x in range(out_w):
            for oc in range(out_c):
                window_flat = []
                kernel_flat = []
                for ic in range(in_c):
                    for dy in range(3):
                        for dx in range(3):
                            window_flat.append(input_map[y + dy][x + dx][ic])
                            kernel_flat.append(weights[oc][ic][dy][dx])
                            
                out_map[y][x][oc] = dsp48_mac(window_flat, kernel_flat, biases[oc], apply_relu)
    return out_map

def padder2d(input_map):
    in_h = len(input_map)
    in_w = len(input_map[0])
    in_c = len(input_map[0][0])
    out_h = in_h + 2
    out_w = in_w + 2
    out_map = [[[0 for _ in range(in_c)] for _ in range(out_w)] for _ in range(out_h)]
    for y in range(in_h):
        for x in range(in_w):
            for c in range(in_c):
                out_map[y + 1][x + 1][c] = input_map[y][x][c]
    return out_map

def max_pool2d_2x2(input_map):
    in_h = len(input_map)
    in_w = len(input_map[0])
    in_c = len(input_map[0][0])
    out_h = in_h // 2
    out_w = in_w // 2
    out_map = [[[0 for _ in range(in_c)] for _ in range(out_w)] for _ in range(out_h)]
    for y in range(out_h):
        for x in range(out_w):
            for c in range(in_c):
                val1 = input_map[y*2][x*2][c]
                val2 = input_map[y*2][x*2 + 1][c]
                val3 = input_map[y*2 + 1][x*2][c]
                val4 = input_map[y*2 + 1][x*2 + 1][c]
                out_map[y][x][c] = max(val1, val2, val3, val4)
    return out_map

def flatten3d(input_map):
    out_flat = []
    for y in range(len(input_map)):
        for x in range(len(input_map[0])):
            for c in range(len(input_map[0][0])):
                out_flat.append(input_map[y][x][c])
    return out_flat

def dense_layer(input_vector, weights, biases, out_c, apply_relu=True):
    """
    Computes Dense output using hardware-aligned transposed weights:
    weights shape: (num_passes, in_size, parallel_macs)
    """
    in_size = len(input_vector)
    parallel_macs = len(weights[0][0]) # Expected to be 8 based on hardware
    
    out_vector = [0] * out_c
    
    for oc in range(out_c):
        # Calculate exactly where this neuron's weights are in the 3D hardware matrix
        pass_idx = oc // parallel_macs
        mac_idx = oc % parallel_macs
        
        # Extract the weight vector for this specific output neuron across all inputs
        weights_flat = [weights[pass_idx][ic][mac_idx] for ic in range(in_size)]
        
        # Run through the bit-exact MAC
        out_vector[oc] = dsp48_mac(input_vector, weights_flat, biases[oc], apply_relu)
        
    return out_vector

# ─────────────────────────────────────────────
#  Phase 1.3: Full Network Forward Pass
# ─────────────────────────────────────────────
def load_q88_weights(layer_name):
    w = np.load(f"export_q88/{layer_name}_weights.npy").tolist()
    b = np.load(f"export_q88/{layer_name}_biases.npy").tolist()
    return w, b

def run_full_inference(img_path):
    print(f"\n--- Starting Full Hardware Simulation for: {img_path} ---")
    
    # 1. Load Image
    img = Image.open(img_path).convert('L').resize((32, 32))
    pixel_data = np.array(img, dtype=int)
    
    # 2. Pi-Side Padding (32x32 -> 34x34)
    padded_img = np.pad(pixel_data, pad_width=1, mode='constant', constant_values=0)
    
    # --- THE FIX: SCALE TO Q8.8 ---
    # Keras trained with 0.0-1.0. Our spi.py divides by 255.0 and multiplies by 256.0.
    # We use np.round() and cast to int to exactly mimic the hardware payload.
    hw_map = [[[int(np.round((padded_img[y][x] / 255.0) * 256.0))] for x in range(34)] for y in range(34)]
    
    # 3. Load All Weights 
    print("Loading exported weights...")
    try:
        c1_w, c1_b = load_q88_weights("conv1")
        c2_w, c2_b = load_q88_weights("conv2")
        c3_w, c3_b = load_q88_weights("conv3")
        d1_w, d1_b = load_q88_weights("dense1")
        d2_w, d2_b = load_q88_weights("dense2")
    except FileNotFoundError:
        print("Error: Could not find weight files. Run export_weights.py first.")
        return None

    # 4. Execute Pipeline
    print("Running Pipeline...")
    x = conv2d_valid(hw_map, c1_w, c1_b, apply_relu=True)
    x = max_pool2d_2x2(x)
    x = padder2d(x)
    x = conv2d_valid(x, c2_w, c2_b, apply_relu=True)
    x = max_pool2d_2x2(x)
    x = padder2d(x)
    x = conv2d_valid(x, c3_w, c3_b, apply_relu=True)
    x = max_pool2d_2x2(x)
    x = flatten3d(x)
    
    # Updated to pass out_c parameter
    x = dense_layer(x, d1_w, d1_b, out_c=32, apply_relu=True)
    logits = dense_layer(x, d2_w, d2_b, out_c=3, apply_relu=False) 
    
    return logits

if __name__ == "__main__":
    # Test a Dog image this time!
    test_image = "Doglets/Human/9_Press_Conference_Press_Conference_9_466.jpg"  # Change to your test image path
    
    if os.path.exists(test_image):
        final_logits = run_full_inference(test_image)
        if final_logits:
            print("\n==================================")
            print(" FPGA SIMULATION COMPLETE")
            print("==================================")
            print(f" Raw Q8.8 Logits : {final_logits}")
            print(f" Predicted Class : {['Cat', 'Dog', 'Human'][np.argmax(final_logits)]}")
            print("==================================")
    else:
        print(f"FAIL: Could not find image at '{test_image}'.")

# ... (Keep all your existing golden_model.py functions above this!)

class GoldenModel:
    def __init__(self):
        self.classes = ["Cat", "Dog", "Human"]

    def predict(self, image_path):
        logits = run_full_inference(image_path)
        if logits is None:
            return "Error"
        # Return the string name of the winning class
        return self.classes[np.argmax(logits)]
\end{lstlisting}

\subsection{cpu.py}
\begin{lstlisting}[language=Python]
import sys
import time
import threading
import numpy as np
import tensorflow as tf
import psutil
from scapy.all import conf, Ether, IP, UDP, sendp

# =========================================================================
# 1. NETWORK & INTERFACE CONFIGURATION
# =========================================================================
INTERFACE = "Realtek PCIe GbE Family Controller" 

PC_MAC = "c8:53:09:fa:62:38" 
FPGA_MAC = "e8:6a:64:e7:e8:30"
PC_IP = "192.168.1.10"
FPGA_IP = "192.168.1.20"

# =========================================================================
# 2. HARDWARE & BENCHMARK DESIGN CONSTANTS
# =========================================================================
FPGA_CLOCK_MHZ = 100.0   # 10 ns clock period
NUM_CLASSES = 3          # 3 target classification layers
RESPONSE_LEN = NUM_CLASSES * 2  # 16-bit elements (6 bytes total)

# Hardcoded true Turbo Boost reference clock for your i7-14700HX P-cores
MANUAL_CPU_FREQ_MHZ = 5200.0

# Fixed compensation factor for OS network stack driver traversal
NETWORK_LATENCY_OFFSET = 0.002 

# =========================================================================
# 3. FILE PATHS
# =========================================================================
INPUT_FILE = "cat_0_input.hex"
EXPECTED_OUTPUT_FILE = "cat_11_logits.hex"
TFLITE_MODEL = "tiny_ai_brain.tflite"


class CnnEthernetTester:
    def __init__(self):
        self.running = True
        self.received_logits_bytes = b""
        
        # Load expected testbench metrics
        self.expected_bytes = self.read_hex_logits(EXPECTED_OUTPUT_FILE)
        print(f"Loaded expected logits: {self.expected_bytes.hex().upper()}")

        # Metric and time tracking
        self.fpga_send_time = 0
        self.fpga_recv_time = 0
        self.fpga_rtt = 0
        self.software_time = 0

        # Initialize TensorFlow Lite inference engine
        self.interpreter = tf.lite.Interpreter(model_path=TFLITE_MODEL)
        self.interpreter.allocate_tensors()

        try:
            self.sock = conf.L2socket(iface=INTERFACE)
        except Exception as e:
            print(f"CRITICAL ERROR Opening Layer 2 Socket: {e}")
            sys.exit(1)

        # Spawn isolated concurrent background receiver thread
        self.listener_thread = threading.Thread(target=self.listen_for_fpga)
        self.listener_thread.daemon = True
        self.listener_thread.start()

    def read_hex_image(self, filepath):
        """Extracts the lower 8 bits of pixel data to match VHDL stream layout."""
        byte_data = bytearray()
        with open(filepath, 'r') as f:
            for line in f:
                line = line.strip()
                if line:
                    byte_data.append(bytes.fromhex(line)[1])
        return bytes(byte_data)

    def read_hex_logits(self, filepath):
        """Reads expected 16-bit logits maintaining full precision window."""
        byte_data = bytearray()
        with open(filepath, 'r') as f:
            for line in f:
                line = line.strip()
                if line:
                    byte_data.extend(bytes.fromhex(line))
        return bytes(byte_data[:RESPONSE_LEN])

    def run_software_inference(self, input_bytes):
        print("\n--- RUNNING SOFTWARE INFERENCE ---")
        input_details = self.interpreter.get_input_details()
        output_details = self.interpreter.get_output_details()
        input_shape = input_details[0]['shape']

        try:
            # Parse 1156 bytes into a 34x34 padded array
            padded_array = np.frombuffer(input_bytes, dtype=np.uint8).reshape((34, 34))
            
            # Crop out the 1-pixel outer border to retrieve the valid 32x32 image matrix
            cropped_array = padded_array[1:33, 1:33]
            
            # Normalize to [0.0, 1.0] float representation matching Q8.8 design logic
            normalized_array = cropped_array.astype(np.float32) / 255.0
            input_array = normalized_array.reshape(input_shape)
            
        except ValueError as e:
            print(f"Matrix parsing or reshape transformation failed: {e}")
            return

        self.interpreter.set_tensor(input_details[0]['index'], input_array)

        # Isolated execution timing window
        start_time = time.perf_counter()
        self.interpreter.invoke()
        end_time = time.perf_counter()

        self.software_time = end_time - start_time
        software_output = self.interpreter.get_tensor(output_details[0]['index'])
        
        print(f"Software Execution Time: {self.software_time:.6f} seconds")
        print(f"Software Output Array:   {software_output.flatten()}")

    def listen_for_fpga(self):
        expected_src_mac = bytes.fromhex(FPGA_MAC.replace(':', ''))
        expected_dst_mac = bytes.fromhex(PC_MAC.replace(':', ''))
        
        print("Listening for FPGA response...")
        while self.running:
            try:
                pkt = self.sock.recv(timeout=1)
                if pkt:
                    raw_bytes = bytes(pkt)
                    
                    if len(raw_bytes) >= 42:
                        actual_dst = raw_bytes[0:6]
                        actual_src = raw_bytes[6:12]
                        
                        if actual_dst == expected_dst_mac and actual_src == expected_src_mac:
                            # Immediate arrival timestamp capture
                            self.fpga_recv_time = time.perf_counter()
                            
                            # Bypass L2/L3 protocol headers (14 Eth + 20 IP + 8 UDP)
                            payload = raw_bytes[42:] 
                            
                            if len(payload) >= RESPONSE_LEN:
                                self.received_logits_bytes = payload[:RESPONSE_LEN]
                                self.compare_and_exit()
            except Exception:
                pass 

    def send_input_image(self):
        input_payload = self.read_hex_image(INPUT_FILE)
        print(f"Loaded {len(input_payload)} correct pixel bytes from {INPUT_FILE}")

        # Compute software benchmark configuration
        self.run_software_inference(input_payload)

        # Construct frame forcing static identifiers and checksums to match VHDL simulation
        frame = Ether(src=PC_MAC, dst=FPGA_MAC) / \
                IP(src=PC_IP, dst=FPGA_IP, id=0, chksum=0) / \
                UDP(sport=4660, dport=17767, chksum=0) / \
                input_payload
        
        print("\n--- RUNNING HARDWARE INFERENCE ---")
        print("Sending UDP CNN Input Data...")
        
        self.fpga_send_time = time.perf_counter()
        sendp(frame, iface=INTERFACE, verbose=False)

    def compare_performance(self):
        print("\n--- PERFORMANCE & CLOCK SCALING ANALYSIS ---")
        
        # Calculate true hardware execution window 
        self.fpga_rtt = self.fpga_recv_time - self.fpga_send_time
        hw_time_adjusted = max(self.fpga_rtt - NETWORK_LATENCY_OFFSET, 0.000001)
        
        print(f"FPGA Round Trip Time (RTT)   : {self.fpga_rtt:.6f} sec")
        print(f"FPGA Est. Time (w/o network) : {hw_time_adjusted:.6f} sec")
        print(f"CPU Time                     : {self.software_time:.6f} sec")
        
        print(f"\nCPU Active Operational Clock : {MANUAL_CPU_FREQ_MHZ:.2f} MHz")
        print(f"FPGA Reference Design Clock  : {FPGA_CLOCK_MHZ:.2f} MHz")

        # Temporal domain converted to absolute core clock cycle counts
        cpu_cycles = self.software_time * (MANUAL_CPU_FREQ_MHZ * 1e6)
        fpga_cycles = hw_time_adjusted * (FPGA_CLOCK_MHZ * 1e6)
        
        print(f"\nCalculated Active CPU Cycles : {int(cpu_cycles):,}")
        print(f"Calculated Running FPGA Cycles: {int(fpga_cycles):,}")
        
        speedup_absolute = self.software_time / hw_time_adjusted
        speedup_scaled = cpu_cycles / fpga_cycles

        print(f"\nAbsolute Speedup (Time-domain)     : {speedup_absolute:.2f}x")
        print(f"Architectural Speedup (Cycle-scale): {speedup_scaled:.2f}x")

    def compare_and_exit(self):
        print("\n--- FUNCTIONAL COMPARISON RESULTS ---")
        print(f"Expected (Hex) : {self.expected_bytes.hex().upper()}")
        print(f"Received (Hex) : {self.received_logits_bytes.hex().upper()}")

        if self.received_logits_bytes == self.expected_bytes:
            print("\n✅ SUCCESS! The FPGA outputs match the testbench perfectly.")
        else:
            print("\n❌ MISMATCH DETECTED!")
        
        self.compare_performance()
        
        # Explicitly clean up Layer 2 socket binding to decrease teardown warning noise
        try:
            self.sock.close()
        except Exception:
            pass
            
        self.running = False

    def wait(self):
        timeout_timer = 5
        while self.running and timeout_timer > 0:
            time.sleep(1)
            timeout_timer -= 1
        
        if timeout_timer == 0 and self.running:
            print("\n[!] TIMEOUT: Did not receive a response from the FPGA in time.")
            self.running = False


if __name__ == "__main__":
    tester = CnnEthernetTester()
    tester.send_input_image()
    tester.wait()
\end{lstlisting}

\subsection{spi.py}
\begin{lstlisting}[language=Python]
import spidev
import time
import numpy as np
from PIL import Image
from gpiozero import Button, DigitalOutputDevice

class FPGACNN:
    # Running at the verified 8 MHz with Dummy CS on GPIO 16
    def __init__(self, bus=0, device=0, speed_hz=8000000, interrupt_pin=25, cs_pin=16):
        self.speed_hz = speed_hz
        self.spi = spidev.SpiDev()
        self.spi.open(bus, device)
        self.spi.max_speed_hz = self.speed_hz
        self.spi.mode = 0
        
        # --- THE DUMMY CS FIX ---
        self.cs_n = DigitalOutputDevice(cs_pin, initial_value=True)
        self.data_ready = Button(interrupt_pin, pull_up=False)

    def preprocess_image(self, image_path):
        """Loads image, zero-pads to 34x34, and converts to Q8.8 byte payload."""
        img = Image.open(image_path).convert('L').resize((32, 32))
        pixel_data = np.array(img, dtype=np.float32)
        
        padded_img = np.pad(pixel_data, pad_width=1, mode='constant', constant_values=0)
        padded_img = (padded_img / 255.0) * 256.0
        fixed_point_pixels = np.round(padded_img).astype(np.int16).flatten().tolist()
        
        spi_payload = []
        for pixel in fixed_point_pixels:
            spi_payload.extend([(pixel >> 8) & 0xFF, pixel & 0xFF])
            
        return spi_payload

    def run_inference(self, image_path):
        payload = self.preprocess_image(image_path)
        
        try:
            self.cs_n.off() # Start Transaction
            
            # --- LAP 1: SPI Transmit ---
            t0 = time.perf_counter()
            self.spi.xfer2(payload)
            
            # --- LAP 2: FPGA Compute ---
            t1 = time.perf_counter() 
            if self.data_ready.wait_for_active(timeout=2.0):
                
                # --- LAP 3: SPI Receive ---
                t2 = time.perf_counter()
                rx_data = self.spi.xfer2([0x00] * 6) 
                t3 = time.perf_counter()
                
                # Calculate the exact millisecond breakdown
                timing = {
                    "spi_tx_ms": (t1 - t0) * 1000.0,
                    "fpga_compute_ms": (t2 - t1) * 1000.0,
                    "spi_rx_ms": (t3 - t2) * 1000.0,
                    "total_system_ms": (t3 - t0) * 1000.0
                }
                
                # Reconstruct 16-bit Signed Q8.8 Logits
                logits = []
                for i in range(0, 6, 2):
                    val = (rx_data[i] << 8) | rx_data[i+1]
                    if val > 32767:
                        val -= 65536 
                    logits.append(val)
                    
                return logits, timing
            else:
                print(f"Timeout waiting for DATA_READY on {image_path}")
                return None, None
                
        finally:
            self.cs_n.on() # End Transaction safely

    def close(self):
        self.spi.close()
\end{lstlisting}

\subsection{benchmark.py}
\begin{lstlisting}[language=Python]
import glob
import numpy as np
import time
import random
from PIL import Image
from gpiozero import Button, DigitalOutputDevice

# Import Custom Hardware & SW Reference
from spi import FPGACNN
from golden_model import run_full_inference

# Import LiteRT (TensorFlow Lite)
try:
    from ai_edge_litert.interpreter import Interpreter
except ImportError:
    print("Error: LiteRT not found. Please install tflite_runtime or ai_edge_litert.")
    exit()

# --- THE LOOPBACK CALIBRATION SETUP ---
# Wire GPIO 17 (Physical Pin 11) directly to GPIO 27 (Physical Pin 13)!
fake_fpga_out = DigitalOutputDevice(17, initial_value=False)
os_listener = Button(27, pull_up=False)

def measure_os_overhead():
    """Forces the Pi to interrupt itself to measure the current OS latency."""
    fake_fpga_out.off()
    time.sleep(0.01) # Let the line settle
    
    t_start = time.perf_counter()
    fake_fpga_out.on() # Trigger the interrupt!
    os_listener.wait_for_active(timeout=1.0)
    t_end = time.perf_counter()
    
    fake_fpga_out.off()
    return (t_end - t_start) * 1000.0 # Return overhead in ms

def prep_tflite_image(image_path):
    """Prepares the image exactly how TFLite expects it: (1, 32, 32, 1) Float32"""
    img = Image.open(image_path).convert('L').resize((32, 32))
    img_array = np.array(img, dtype=np.float32)
    return np.expand_dims(img_array, axis=(0, -1))

def main():
    print("======================================================")
    print("  HARDWARE vs SOFTWARE PERFORMANCE BENCHMARK SUITE  ")
    print("======================================================")
    
    # 1. Initialize Hardware Engine
    fpga_engine = FPGACNN(speed_hz=8000000, interrupt_pin=25, cs_pin=16)
    
    # 2. Initialize TFLite Software Engine
    interpreter = Interpreter(model_path="tiny_ai_brain.tflite")
    interpreter.allocate_tensors()
    in_idx = interpreter.get_input_details()[0]['index']
    out_idx = interpreter.get_output_details()[0]['index']
    
    # --- WARM UP THE CPU ---
    print("Warming up the CPU cache...")
    dummy_input = np.zeros((1, 32, 32, 1), dtype=np.float32)
    interpreter.set_tensor(in_idx, dummy_input)
    interpreter.invoke()
    
    # 3. Find and Randomize Test Images
    test_images = glob.glob("Doglets/*/*.jpg")
    if not test_images:
        print("Error: Could not find test images in 'Doglets/*/*.jpg'")
        return
        
    random.shuffle(test_images) # THE FIX: Randomize the testing order!
    
    matches = 0
    total_images = min(10, len(test_images)) 
    classes = ["Cat", "Dog", "Human"]
    
    for img_path in test_images[:total_images]: 
        print(f"\nEvaluating: {img_path}")
        
        # ----------------------------------------------------
        # TEST A: THE SOFTWARE BENCHMARK (TFLite)
        # ----------------------------------------------------
        sw_input = prep_tflite_image(img_path)
        t0_sw = time.perf_counter()
        interpreter.set_tensor(in_idx, sw_input)
        interpreter.invoke()
        sw_raw = interpreter.get_tensor(out_idx)[0]
        sw_time_ms = (time.perf_counter() - t0_sw) * 1000.0
        sw_pred = classes[np.argmax(sw_raw)]
        
        # ----------------------------------------------------
        # DYNAMIC OS CALIBRATION
        # ----------------------------------------------------
        current_os_overhead = measure_os_overhead()
        
        # ----------------------------------------------------
        # TEST B: THE HARDWARE BENCHMARK (FPGA)
        # ----------------------------------------------------
        fpga_logits, hw_timing = fpga_engine.run_inference(img_path)
        if fpga_logits is None:
            continue
        fpga_pred = classes[np.argmax(fpga_logits)]
        
        # ----------------------------------------------------
        # TEST C: THE ACCURACY BENCHMARK (Golden Model)
        # ----------------------------------------------------
        golden_logits = run_full_inference(img_path)
        is_exact = "YES" if fpga_logits == golden_logits else "NO"
        if is_exact == "YES":
            matches += 1

        # ----------------------------------------------------
        # CALCULATE TRUE FPGA MATH TIME
        # ----------------------------------------------------
        raw_measured = hw_timing['fpga_compute_ms']
        true_fpga_math = raw_measured - current_os_overhead
        
        # Prevent "negative time" caused by Linux Scheduler Jitter
        if true_fpga_math < 0.001:
            true_fpga_math = 0.001 

        # ----------------------------------------------------
        # PRINT THE RESULTS TABLE
        # ----------------------------------------------------
        print(f"| Metric                   | Pi 5 SW (TFLite) | FPGA HW (Artix-7) |")
        print(f"|--------------------------|------------------|-------------------|")
        print(f"| Prediction               | {sw_pred:<16} | {fpga_pred:<17} |")
        print(f"| 100% Bit-Exact Parity?   | N/A              | {is_exact:<17} |")
        print(f"| Compute Time (Pure Math) | {sw_time_ms:>13.3f} ms | {true_fpga_math:>14.3f} ms |")
        print(f"| SPI Transfer Time (I/O)  | N/A              | {hw_timing['spi_tx_ms'] + hw_timing['spi_rx_ms']:>14.3f} ms |")
        print(f"| Linux OS Overhead Delay  | N/A              | {current_os_overhead:>14.3f} ms |") 
        print(f"| Total System Latency     | {sw_time_ms:>13.3f} ms | {hw_timing['total_system_ms']:>14.3f} ms |")
        print(f"|--------------------------|------------------|-------------------|")
        
    print("\n======================================================")
    print(f" BENCHMARK COMPLETE ")
    print(f" Verified Math Parity: {matches}/{total_images} ({(matches/total_images)*100}%)")
    print("======================================================")
    
    fpga_engine.close()

if __name__ == "__main__":
    main()
\end{lstlisting}

\subsection{generate_testbench_data.py}
\begin{lstlisting}[language=Python]
import os
import numpy as np
from golden_model import (
    run_full_inference, dsp48_mac, conv2d_valid, 
    padder2d, max_pool2d_2x2, flatten3d, dense_layer, load_q88_weights
)
from PIL import Image

# ─────────────────────────────────────────────
#  Phase 1.4: Testbench Data Generator
# ─────────────────────────────────────────────

OUTPUT_DIR = "testbench_data"

def to_hex_16bit(val):
    """Converts a Q8.8 integer to a 4-char hex string."""
    return f"{int(val) & 0xFFFF:04X}"

def save_map_to_hex(filename, data_map):
    """Flattens a 3D or 1D map and saves as a .hex file for VHDL $readmemh."""
    # Flatten if it's a 3D map (H, W, C)
    if isinstance(data_map[0], list) and isinstance(data_map[0][0], list):
        flat_list = []
        for y in range(len(data_map)):
            for x in range(len(data_map[0])):
                for c in range(len(data_map[0][0])):
                    flat_list.append(data_map[y][x][c])
    else:
        flat_list = data_map # Already 1D (like logits or flattened layer)

    with open(os.path.join(OUTPUT_DIR, filename), 'w') as f:
        for val in flat_list:
            f.write(to_hex_16bit(val) + "\n")

def generate_full_trace(img_path, label):
    print(f"\nTracing layers for: {label} ({img_path})")
    
    # 1. Pre-process Image (Master Plan Sec 7.3)
    img = Image.open(img_path).convert('L').resize((32, 32))
    pixel_data = np.array(img, dtype=int)
    padded_img = np.pad(pixel_data, pad_width=1, mode='constant', constant_values=0)
    hw_map = [[[int(padded_img[y][x])] for x in range(34)] for y in range(34)]
    
    # Save input image hex (34x34 = 1156 words)
    save_map_to_hex(f"{label}_0_input.hex", hw_map)

    # 2. Load Weights
    c1_w, c1_b = load_q88_weights("conv1")
    c2_w, c2_b = load_q88_weights("conv2")
    c3_w, c3_b = load_q88_weights("conv3")
    d1_w, d1_b = load_q88_weights("dense1")
    d2_w, d2_b = load_q88_weights("dense2")

    # 3. Run and Save Every Layer Step
    print("  Processing Conv1...")
    x = conv2d_valid(hw_map, c1_w, c1_b, apply_relu=True)
    save_map_to_hex(f"{label}_1_conv1.hex", x)
    
    print("  Processing Pool1...")
    x = max_pool2d_2x2(x)
    save_map_to_hex(f"{label}_2_pool1.hex", x)
    
    print("  Processing Pad1...")
    x = padder2d(x)
    save_map_to_hex(f"{label}_3_pad1.hex", x)
    
    print("  Processing Conv2...")
    x = conv2d_valid(x, c2_w, c2_b, apply_relu=True)
    save_map_to_hex(f"{label}_4_conv2.hex", x)
    
    print("  Processing Pool2...")
    x = max_pool2d_2x2(x)
    save_map_to_hex(f"{label}_5_pool2.hex", x)
    
    print("  Processing Pad2...")
    x = padder2d(x)
    save_map_to_hex(f"{label}_6_pad2.hex", x)
    
    print("  Processing Conv3...")
    x = conv2d_valid(x, c3_w, c3_b, apply_relu=True)
    save_map_to_hex(f"{label}_7_conv3.hex", x)
    
    print("  Processing Pool3...")
    x = max_pool2d_2x2(x)
    save_map_to_hex(f"{label}_8_pool3.hex", x)
    
    print("  Processing Flatten...")
    x = flatten3d(x)
    save_map_to_hex(f"{label}_9_flatten.hex", x)
    
    print("  Processing Dense1...")
    # UPDATED: Added out_c=32
    x = dense_layer(x, d1_w, d1_b, out_c=32, apply_relu=True)
    save_map_to_hex(f"{label}_10_dense1.hex", x)
    
    print("  Processing Dense2...")
    # UPDATED: Added out_c=3
    logits = dense_layer(x, d2_w, d2_b, out_c=3, apply_relu=False)
    save_map_to_hex(f"{label}_11_logits.hex", logits)
    
    print(f"Done. Files saved to {OUTPUT_DIR}/")

if __name__ == "__main__":
    if not os.path.exists(OUTPUT_DIR):
        os.makedirs(OUTPUT_DIR)

    # We use 3 images for the answer key (one per class)
    test_images = {
        "cat": "Doglets/Cat/0.jpg",
        "dog": "Doglets/Dog/0.jpg",
        "human": "Doglets/Human/image_025556.png" # Update with a real path from your Human folder
    }

    for label, path in test_images.items():
        if os.path.exists(path):
            generate_full_trace(path, label)
        else:
            print(f"Warning: Could not find {path}, skipping {label} trace.")
\end{lstlisting}

\subsection{convert_h5_to_tflite.py}
\begin{lstlisting}[language=Python]
import tensorflow as tf
import sys
import os

def convert_h5_to_tflite(h5_path):
    if not os.path.exists(h5_path):
        print(f"Error: {h5_path} does not exist.")
        return

    # Create tflite path in the same directory with same base name
    base_name = os.path.splitext(h5_path)[0]
    tflite_path = base_name + ".tflite"

    try:
        print(f"Loading model from {h5_path}...")
        model = tf.keras.models.load_model(h5_path)
        
        print("Converting model to TFLite...")
        converter = tf.lite.TFLiteConverter.from_keras_model(model)
        
        # Uncomment the following line if you want to apply post-training quantization
        # converter.optimizations = [tf.lite.Optimize.DEFAULT]
        
        tflite_model = converter.convert()
        
        print(f"Saving TFLite model to {tflite_path}...")
        with open(tflite_path, 'wb') as f:
            f.write(tflite_model)
            
        print("Conversion completed successfully!")
        
    except Exception as e:
        print(f"An error occurred during conversion: {e}")

if __name__ == "__main__":
    # Use the specific .h5 file found in the directory by default
    # Change this if you have a different filename
    default_h5 = "tiny_ai_brain.h5"
    
    # Allow passing filename as argument, e.g., `python convert_h5_to_tflite.py my_model.h5`
    if len(sys.argv) > 1:
        h5_file = sys.argv[1]
    else:
        h5_file = default_h5
        
    convert_h5_to_tflite(h5_file)

\end{lstlisting}


\end{document}
